%% This is a LaTeX document.  Hey, Emacs, -*- latex -*- , get it?
\documentclass[12pt,a4paper]{article}
\usepackage[a4paper,margin=2.5cm]{geometry}
\usepackage[shorthands=off,french]{babel}
\usepackage[utf8]{inputenc}
\usepackage[T1]{fontenc}
\usepackage{lmodern}
\usepackage{newtxtext}
% A tribute to the worthy AMS:
\usepackage{amsmath}
\usepackage{amsfonts}
\usepackage{amssymb}
\usepackage{amsthm}
%
\usepackage{mathrsfs}
\usepackage{wasysym}
\usepackage{url}
\usepackage{mathpartir}
\usepackage{flagderiv}
%
\usepackage{makeidx}
%% Self-note: compile index with:
%% texindy -C utf8 -L french notes-mitro206.idx
%
\usepackage{graphicx}
%\usepackage{tikz}
%\usetikzlibrary{arrows,automata,calc}
\usepackage[usenames,dvipsnames]{xcolor}
\usepackage{tikz}
\usetikzlibrary{arrows}
\usepackage[hyperindex=false]{hyperref}
%
\theoremstyle{definition}
\newtheorem{comcnt}{Tout}[subsection]
\newcommand\thingy{%
\refstepcounter{comcnt}\medskip\noindent\textbf{\thecomcnt.} }
\newcommand\exercice[1][Exercice]{%
\refstepcounter{comcnt}\bigskip\noindent\textbf{#1~\thecomcnt.}\par\nobreak}
\newtheorem{defn}[comcnt]{Définition}
\newtheorem{prop}[comcnt]{Proposition}
\newtheorem{lem}[comcnt]{Lemme}
\newtheorem{thm}[comcnt]{Théorème}
\newtheorem{cor}[comcnt]{Corollaire}
\newtheorem{scho}[comcnt]{Scholie}
\newtheorem{algo}[comcnt]{Algorithme}
\renewcommand{\qedsymbol}{\smiley}
% Poor man's chapter commands:
% Chapitres dans la classe article.
\newcounter{chapter}
\makeatletter
\newcommand*\l@chapter[2]{%
  \addpenalty{-\@highpenalty}%
  \addvspace{1.0em plus\p@}%
  \setlength\@tempdima{3em}%
  \begingroup
    \parindent\z@ \rightskip\@pnumwidth
    \parfillskip-\@pnumwidth
    \leavevmode\bfseries
    \advance\leftskip\@tempdima
    \hskip-\leftskip
    #1\nobreak\hfil\nobreak\hb@xt@\@pnumwidth{\hss #2}\par
  \endgroup
}
\makeatother
\newcommand{\chapter}[1]{%
  \par\bigskip
  \refstepcounter{chapter}%
  \addcontentsline{toc}{chapter}{\Roman{chapter}.~#1}%
  \begin{center}
    \Large\bfseries
    \Roman{chapter}.~#1\par
  \end{center}
  \medskip
}
%
\newcommand{\review}[1]{\textcolor{red}{#1}}
%
\newcommand{\dbllangle}{\mathopen{\langle\!\langle}}
\newcommand{\dblrangle}{\mathclose{\rangle\!\rangle}}
\newcommand{\dottedlimp}{\mathbin{\dot\Rightarrow}}
\newcommand{\dottedland}{\mathbin{\dot\land}}
\newcommand{\dottedlor}{\mathbin{\dot\lor}}
\newcommand{\dottedtop}{\mathord{\dot\top}}
\newcommand{\dottedbot}{\mathord{\dot\bot}}
\newcommand{\dottedneg}{\mathop{\dot\neg}}
\newcommand{\id}{\operatorname{id}}
\newcommand{\limp}{\Longrightarrow}
\newcommand{\cps}{\mathrm{CPS}}
\mathchardef\emdash="07C\relax
%
\DeclareUnicodeCharacter{00A0}{~}
\DeclareUnicodeCharacter{2026}{...}
%
%
\newcommand{\spaceout}{\hskip1emplus2emminus.5em}
\newif\ifcorrige
\corrigetrue
\newenvironment{corrige}%
{\ifcorrige\relax\else\setbox0=\vbox\bgroup\fi%
\smallbreak\footnotesize\noindent{\underbar{\textit{Corrigé.}}\quad}}
{{\hbox{}\nobreak\hfill\checkmark}%
\ifcorrige\par\smallbreak\else\egroup\par\fi}
%
\newcommand{\defin}[2][]{\def\latexsucks{#1}\ifx\latexsucks\empty\index{#2}\else\index{\latexsucks}\fi\textbf{#2}}
%
%
%
\begin{document}
\title{CSC-3C34-TP / INF110\\
(Logique et Fondements de l'Informatique)
\\\footnotesize Notes de cours provisoires}
\author{David A. Madore\\
{\footnotesize Télécom Paris}\\
\texttt{david.madore@}(nom de l'école)\texttt{.fr}}
\date{2023–2027}
\maketitle
{\footnotesize\center{\url{http://perso.enst.fr/madore/inf110/notes-inf110.pdf}}\par}
{\tiny
\immediate\write18{sh ./vc > vcline.tex}
\begin{center}
Git: \input{vcline.tex}
\end{center}
\immediate\write18{echo ' (stale)' >> vcline.tex}
\par}

\section{Calculabilité}\label{section-computability}

\begin{center}
\textcolor{red}{\textbf{Premier brouillon produit avec l'aide d'un
modèle d'IA.  Ce texte doit impérativement être relu, corrigé et
vérifié par l'enseignant avant toute diffusion.}}
\end{center}

\subsection{Introduction}\label{subsection-computability-introduction}

\thingy La \defin{calculabilité} se situe à l'interface entre la logique
mathématique et l'informatique théorique.  Elle est née de
préoccupations venues de la logique, notamment des travaux de Hilbert
et de Gödel, et elle a conduit à certains des premiers concepts
fondamentaux de l'informatique théorique, comme le $\lambda$-calcul et
la machine de Turing.

Son but est d'étudier les limites de ce que peut ou ne peut pas faire
un algorithme.  Contrairement à la théorie de la complexité, elle ne
cherche pas à mesurer précisément le temps ou la mémoire nécessaires à
un calcul.  Elle autorise des ressources arbitrairement grandes, à
condition qu'elles restent finies pour chaque calcul qui termine.  Elle
permet ainsi de démontrer des résultats négatifs de la forme
«~\emph{aucun} algorithme ne peut résoudre ce problème~».

On peut également considérer des notions relatives de calculabilité,
dans lesquelles un algorithme dispose d'informations supplémentaires
sous la forme d'un oracle, ainsi que diverses généralisations dites de
calculabilité supérieure.  Ces développements ne seront toutefois pas
au centre de cette partie.

\thingy Quelques noms permettent de situer historiquement le sujet.
Le mot «~algorithme~» dérive du nom de Muḥammad ibn Mūsá
al-\b{H}wārizmī (vers 780--vers 850).  Blaise Pascal
(1623--1662) construisit une machine à calculer, tandis que Charles
Babbage (1791--1871) conçut son \textit{Analytical Engine}, machine
programmable qui, dans une idéalisation convenable, aurait eu la
puissance de calcul d'une machine de Turing.  Ada, née Byron, comtesse
de Lovelace (1815--1852), est notamment associée aux premiers travaux
de programmation pour cette machine.

Du côté mathématique, Richard Dedekind (1831--1916) introduisit des
définitions par récursion primitive.  David Hilbert (1862--1943)
formula l'\textit{Entscheidungsproblem}, qui demandait essentiellement
une procédure générale permettant de décider la vérité des énoncés
mathématiques d'une certaine forme.  Jacques Herbrand (1908--1931),
Kurt Gödel (1906--1978) et Stephen C. Kleene (1909--1994)
contribuèrent à la théorie des fonctions récursives.  Alonzo Church
(1903--1995) introduisit le $\lambda$-calcul, tandis qu'Alan M. Turing
(1912--1954) définit les machines qui portent son nom et démontra
l'indécidabilité du problème de l'arrêt.  Emil Post (1897--1954)
étudia notamment les ensembles calculablement énumérables.  Enfin,
Haskell Curry (1900--1982) développa la logique combinatoire et
contribua à mettre en évidence les liens entre preuves et typage.

\thingy Une première tentative, encore informelle, consiste à dire
qu'une fonction $f$ est \defin{calculable} lorsqu'il existe un
algorithme qui, prenant en entrée un élément $x$ du domaine de
définition de $f$, termine en temps fini et renvoie la valeur $f(x)$.

Cette formulation fait immédiatement apparaître plusieurs difficultés.
Que faut-il entendre exactement par «~algorithme~»~?  Quelles sortes
de valeurs un algorithme peut-il accepter ou renvoyer~?  Comment
décrire les situations dans lesquelles il ne termine pas~?  Enfin, il
faut distinguer le programme lui-même, ou son \emph{intention}, de la
fonction qu'il calcule, ou son \emph{extension}~: plusieurs programmes
différents peuvent calculer exactement la même fonction.

\thingy La calculabilité ne s'intéresse pas à l'efficacité des
algorithmes, mais seulement à leur terminaison en temps fini.  Par
exemple, pour déterminer si un entier $n$ est premier, on peut tester
tous les produits $i\times j$ avec $2\leq i,j\leq n-1$.  Cet
algorithme est extrêmement inefficace, mais il termine pour toute
entrée $n$ et suffit donc du point de vue de la calculabilité.

De même, la calculabilité n'a pas peur des grands entiers.  Considérons
la \defin{fonction d'Ackermann} définie par
\[
\begin{aligned}
A(m,n,0) &= m+n,\\
A(m,1,k+1) &= m,\\
A(m,n+1,k+1) &= A(m,\,A(m,n,k+1),\,k).
\end{aligned}
\]
Cette définition décrit un algorithme par appels récursifs, et ces
appels terminent.  La fonction est donc calculable.  Ses valeurs
deviennent néanmoins gigantesques très rapidement~:
\[
A(2,6,3)
 = 2^{2^{2^{2^{2^2}}}}
 = 2^{2^{65\,536}},
\]
et $A(2,4,4)=A(2,65\,536,3)$ est déjà inimaginablement grand.
Une valeur telle que $A(100,100,100)$ est tout à fait hors de portée
d'un ordinateur réel, sans que cela empêche la fonction $A$ d'être
calculable.

\thingy On peut aborder la notion de calculabilité de plusieurs
manières.

L'approche informelle considère un algorithme comme un calcul
\emph{finitiste}, mené par un humain ou une machine selon des
instructions précises, en temps fini et sur des données finies.

L'approche pragmatique identifie les fonctions calculables à celles
que l'on peut programmer dans un langage général dit
«~Turing-complet~», par exemple Python, Java, C ou OCaml, à condition
d'idéaliser ce langage.  On suppose notamment que les entiers n'ont
aucune limite de taille et que la mémoire disponible n'est pas bornée
à l'avance.  On écarte également les véritables sources de hasard et
de non-déterminisme.

Enfin, plusieurs approches formelles ont été proposées~:
\begin{itemize}
\item les fonctions générales récursives de Herbrand, Gödel et Kleene~;
\item le $\lambda$-calcul de Church, proche des langages fonctionnels~;
\item les machines de Turing~;
\item les machines à registres, notamment étudiées par Post.
\end{itemize}

\thingy Le résultat mathématique fondamental est que les principales
formalisations précédentes conduisent à la même classe de fonctions.

\begin{thm}[Post, Turing]\label{church-turing-equivalence}
Les classes de fonctions partielles
$\mathbb{N}\dasharrow\mathbb{N}$ suivantes coïncident~:
\begin{enumerate}
\item les fonctions générales récursives~;
\item les fonctions représentables dans le $\lambda$-calcul~;
\item les fonctions calculables par une machine de Turing.
\end{enumerate}
\end{thm}

On parle donc de \defin{calculabilité au sens de Church-Turing}.
Tous les langages de programmation généraux usuels, convenablement
idéalisés, calculent eux aussi exactement ces fonctions~: c'est ce que
signifie, du point de vue de la calculabilité, l'expression
«~Turing-complet~».

Il faut distinguer ce théorème de deux affirmations plus générales.  La
\defin[Church-Turing (thèse de)]{thèse de Church-Turing}, de nature
philosophique, affirme que cette définition formelle saisit exactement
la notion intuitive d'algorithme finitiste.  Une conjecture physique
plus forte affirme que les calculs réalisables mécaniquement dans
l'Univers, avec un temps et une énergie finis mais sans borne fixée à
l'avance, sont précisément ceux de Church-Turing.  Cette conjecture
reste la même si l'on autorise un ordinateur quantique~: un ordinateur
quantique pourrait apporter des gains de complexité, mais ne devrait
pas calculer davantage de fonctions au sens considéré ici.

\thingy Nous allons étudier trois descriptions de ces mêmes fonctions
calculables.

Les \defin{fonctions générales récursives} sont particulièrement
commodes du point de vue mathématique.  Toutes les données y sont
représentées par des entiers, les fonctions considérées étant de la
forme
\[
\mathbb{N}^k\dasharrow\mathbb{N}.
\]
Leur définition est inductive et permet d'associer des numéros aux
programmes.

Les \defin{machines de Turing} constituent une modélisation très simple
d'un ordinateur.  Elles travaillent sur une bande illimitée \textit{a
priori}, qui joue le rôle de mémoire.  Leur caractère algorithmique est
particulièrement manifeste.  Elles sont aussi commodes pour étudier la
complexité, même si cette question ne sera pas notre préoccupation
principale ici.

Enfin, le \defin{$\lambda$-calcul} pur non typé est un système
symbolique proche des langages de programmation fonctionnels comme
Lisp, Haskell ou OCaml.  Il est relativement facile d'y écrire de
véritables programmes, mais son étude fait intervenir plusieurs
subtilités liées notamment aux variables et à la stratégie
d'évaluation.

\subsection{Données finies, codages et fonctions partielles}
\label{subsection-finite-data-and-partial-functions}

\thingy Un algorithme travaille sur des \defin{données finies}.  Cette
expression recouvre les booléens, les entiers, les chaînes de
caractères, les listes, les tableaux et les structures finies plus
complexes, par exemple les graphes finis.

Deux attitudes sont possibles.  Le \defin{typage} distingue les
différentes sortes de données et précise les opérations permises sur
chacune d'elles.  Cette approche est élégante et proche de
l'informatique réelle.  À l'opposé, le \defin[Gödel
(codage de)]{codage de Gödel} représente toutes les données finies par
des entiers naturels.  Cette seconde approche est moins élégante mais
simplifie considérablement la définition de la calculabilité, puisqu'il
suffit alors d'étudier des fonctions
$\mathbb{N}^k\dasharrow\mathbb{N}$, voire seulement
$\mathbb{N}\dasharrow\mathbb{N}$.

\thingy\label{godel-coding} Pour coder un couple d'entiers, on peut par
exemple poser
\[
\langle m,n\rangle
 := m+\frac{1}{2}(m+n)(m+n+1).
\]
Cette formule définit une bijection calculable
$\mathbb{N}^2\to\mathbb{N}$, dont la réciproque est elle aussi
calculable.  On peut donc aussi bien manipuler un couple $(m,n)$ que
l'entier $\langle m,n\rangle$ qui le représente.

On peut ensuite coder une liste finie par
\[
\dbllangle a_0,\ldots,a_{k-1}\dblrangle
 := \langle a_0,
       \langle a_1,
       \langle\cdots,
       \langle a_{k-1},0\rangle+1
       \cdots\rangle+1\rangle+1,
\]
avec la convention $\dbllangle\dblrangle:=0$.  On obtient ainsi une
bijection calculable entre l'ensemble des suites finies d'entiers
naturels et $\mathbb{N}$.

Des codages analogues permettront de représenter des expressions, des
arbres, des configurations de machines et même des programmes par des
entiers.  Le choix précis du codage est sans importance pour la théorie,
pourvu que les opérations d'encodage et de décodage nécessaires soient
calculables.

{\footnotesize
\thingy \emph{Attention~:} ces codages sont des outils théoriques.  Ils
ne constituent généralement pas de bonnes représentations pour un
programme réel, en raison notamment de la taille considérable des
entiers obtenus.
\par}

\thingy Même si l'on cherche principalement à décrire les algorithmes
qui terminent, la définition générale de la calculabilité doit prendre
en compte ceux qui ne terminent pas.  On est ainsi conduit à travailler
avec des \defin{fonctions partielles}.

Une fonction partielle
\[
f\colon\mathbb{N}\dasharrow\mathbb{N}
\]
est une fonction $D\to\mathbb{N}$ définie sur une partie
$D\subseteq\mathbb{N}$.  On écrit
\[
f(n)\downarrow
\]
lorsque $n\in D$, c'est-à-dire lorsque $f(n)$ est définie, et
\[
f(n)\uparrow
\]
lorsque $n\notin D$.

L'expression
\[
f(n)\downarrow=g(m)
\]
signifie que $f(n)$ et $g(m)$ sont toutes deux définies et ont la même
valeur.  Par convention, l'égalité
\[
f(n)=g(m)
\]
entre deux expressions éventuellement partielles signifie qu'elles
sont définies simultanément et qu'elles ont alors la même valeur.
Certains auteurs écrivent plutôt $f(n)\simeq g(m)$ pour cette relation.

Enfin, si l'un des arguments $g_i(\underline{x})$ n'est pas défini, on
convient que
\[
h(g_1(\underline{x}),\ldots,g_k(\underline{x}))\uparrow.
\]
Une fonction partielle définie sur tout $\mathbb{N}^k$ est dite
\defin[totale (fonction)]{totale}.  Une fonction totale est donc un cas
particulier de fonction partielle.

\begin{defn}\label{definition-computable-partial-function}
Une fonction partielle
$f\colon\mathbb{N}^k\dasharrow\mathbb{N}$ est
\defin[calculable (fonction)]{calculable} lorsqu'il existe un algorithme
qui, pour toute entrée $\underline{x}\in\mathbb{N}^k$,
\begin{itemize}
\item termine en temps fini et renvoie $f(\underline{x})$ lorsque
  $f(\underline{x})\downarrow$~;
\item ne termine pas lorsque $f(\underline{x})\uparrow$.
\end{itemize}
\end{defn}

\begin{defn}\label{definition-decidable-semidecidable}
Une partie $A\subseteq\mathbb{N}^k$ est dite
\defin[décidable (ensemble)]{décidable} lorsque sa fonction indicatrice
\[
\mathbf{1}_A\colon\underline{x}\mapsto
\begin{cases}
1&\text{si $\underline{x}\in A$,}\\
0&\text{si $\underline{x}\notin A$}
\end{cases}
\]
est calculable.

Elle est dite \defin[semi-décidable
(ensemble)]{semi-décidable} lorsque la fonction partielle
\[
\underline{x}\mapsto
\begin{cases}
1&\text{si $\underline{x}\in A$,}\\
\uparrow&\text{si $\underline{x}\notin A$}
\end{cases}
\]
est calculable.
\end{defn}

Ainsi, un algorithme de décision répond toujours «~oui~» ou «~non~».
Un algorithme de semi-décision répond «~oui~» lorsque l'entrée
appartient à $A$, mais peut calculer indéfiniment lorsqu'elle n'y
appartient pas.

\thingy Le mot «~récursif~» possède plusieurs sens apparentés mais
distincts.  Il peut signifier «~défini par récurrence~», comme dans les
expressions «~fonction primitive récursive~» et «~fonction générale
récursive~».  Il est aussi employé comme synonyme historique de
«~calculable~».  Enfin, en informatique, un appel est dit récursif
lorsqu'une fonction fait appel à elle-même.  Dans la suite, nous
réserverons autant que possible l'expression «~appel récursif~» à ce
dernier sens.

\subsection{Fonctions primitives récursives}
\label{subsection-primitive-recursive-functions}

\thingy Avant de définir les fonctions générales récursives, qui
coïncideront avec les fonctions calculables, commençons par une classe
plus restreinte~: les \defin[primitive récursive
(fonction)]{fonctions primitives récursives}, abrégées en
\textbf{p.r.}

Cette classe est historiquement antérieure à la définition de la
calculabilité au sens de Church-Turing.  Elle est utile comme première
approximation et correspond informatiquement à des programmes dont
toutes les boucles sont bornées \textit{a priori}.  Une très grande
partie des algorithmes usuels définit des fonctions primitives
récursives, mais ce n'est pas le cas de tous les algorithmes
terminants~: la fonction d'Ackermann fournira un contre-exemple.

\begin{defn}\label{primitive-recursive-definition}
La classe $\mathbf{PR}$ des fonctions primitives récursives est la plus
petite classe de fonctions
$\mathbb{N}^k\dasharrow\mathbb{N}$, pour $k$ variable, qui satisfait
les propriétés suivantes.
\begin{itemize}
\item Elle contient les projections
\[
(x_1,\ldots,x_k)\mapsto x_i.
\]
\item Elle contient toutes les fonctions constantes
\[
\underline{x}\mapsto c.
\]
\item Elle contient la fonction successeur $x\mapsto x+1$.
\item Elle est stable par composition~: si
\[
g_1,\ldots,g_\ell\colon\mathbb{N}^k\dasharrow\mathbb{N}
\quad\text{et}\quad
h\colon\mathbb{N}^\ell\dasharrow\mathbb{N}
\]
sont primitives récursives, alors
\[
\underline{x}\mapsto
h(g_1(\underline{x}),\ldots,g_\ell(\underline{x}))
\]
est primitive récursive.
\item Elle est stable par \defin{récursion primitive}~: si
\[
g\colon\mathbb{N}^k\dasharrow\mathbb{N}
\quad\text{et}\quad
h\colon\mathbb{N}^{k+2}\dasharrow\mathbb{N}
\]
sont primitives récursives, alors la fonction
$f\colon\mathbb{N}^{k+1}\dasharrow\mathbb{N}$ définie par
\[
\begin{aligned}
f(\underline{x},0) &= g(\underline{x}),\\
f(\underline{x},z+1)
 &=h(\underline{x},f(\underline{x},z),z)
\end{aligned}
\]
est primitive récursive.
\end{itemize}
\end{defn}

Toutes les fonctions primitives récursives sont en fait totales.  La
notation $\dasharrow$ reste néanmoins commode, notamment lorsque nous
coderons les programmes~: certains entiers ne coderont aucun programme
bien formé et correspondront alors à une fonction non définie.

\thingy L'addition est primitive récursive, car la fonction
$f(x,z)=x+z$ vérifie
\[
\begin{aligned}
f(x,0)&=x,\\
f(x,z+1)&=f(x,z)+1.
\end{aligned}
\]
La fonction initiale $x\mapsto x$ est une projection, et la fonction
$(x,y,z)\mapsto y+1$ s'obtient par composition de la projection sur
$y$ et du successeur.

La multiplication est ensuite primitive récursive, puisque
$f(x,z)=x\cdot z$ vérifie
\[
\begin{aligned}
f(x,0)&=0,\\
f(x,z+1)&=f(x,z)+x.
\end{aligned}
\]
En utilisant de nouveau une récursion primitive, on obtient que
$(x,z)\mapsto x^z$ est primitive récursive.

La fonction définie par
\[
f(x,y,z)=
\begin{cases}
x&\text{si $z=0$,}\\
y&\text{si $z\geq1$}
\end{cases}
\]
est elle aussi primitive récursive, car
\[
\begin{aligned}
f(x,y,0)&=x,\\
f(x,y,z+1)&=y.
\end{aligned}
\]

Enfin, la fonction prédécesseur tronqué
\[
u\mapsto\max(u-1,0)
\]
est primitive récursive.  On peut en déduire que le sont également la
soustraction tronquée
\[
(u,v)\mapsto\max(u-v,0),
\]
le reste $(u,v)\mapsto u\mathbin{\%}v$ et le quotient entier
$(u,v)\mapsto\lfloor u/v\rfloor$.

\review{Les diapositives présentent la construction explicite du
prédécesseur tronqué et des opérations de division comme exercice.
La démonstration détaillée pourra être ajoutée lors de l'intégration
des feuilles d'exercices.}

\thingy Une autre façon de comprendre la classe $\mathbf{PR}$ consiste
à employer un langage de programmation idéalisé dans lequel~:
\begin{itemize}
\item les variables contiennent des entiers naturels de taille
  arbitraire~;
\item les constantes, affectations, tests d'égalité et conditionnelles
  sont disponibles~;
\item les opérations arithmétiques élémentaires sont disponibles~;
\item les appels de fonctions sont permis, mais pas les appels
  récursifs~;
\item toutes les boucles ont un nombre d'itérations borné
  \textit{a priori}.
\end{itemize}
Les programmes écrits dans un tel langage terminent par construction.
Les fonctions qu'ils calculent sont exactement les fonctions
primitives récursives.

Les fonctions de codage et de décodage des couples introduites
en~\ref{godel-coding}, à savoir
\[
(m,n)\mapsto\langle m,n\rangle,\qquad
\langle m,n\rangle\mapsto m,\qquad
\langle m,n\rangle\mapsto n,
\]
sont elles aussi primitives récursives.

\thingy La classe $\mathbf{PR}$ peut aussi être interprétée comme une
classe de complexité extrêmement large.  En anticipant sur la
définition des machines de Turing, la fonction qui, à une machine $M$
et une configuration $C$, associe la configuration suivante $C'$ est
primitive récursive.  Par récursion primitive, la fonction
\[
(n,M,C)\mapsto C^{(n)}
\]
qui associe à $M$, $C$ et $n$ la configuration atteinte après $n$
étapes est donc primitive récursive.

Il s'ensuit qu'une fonction calculée par une machine de Turing en un
nombre d'étapes borné par une fonction primitive récursive est
elle-même primitive récursive.  Réciproquement, toute fonction
primitive récursive est calculable avec une borne primitive récursive
sur sa complexité.

\review{L'équivalence précise entre «~fonction primitive récursive~»
et «~fonction de complexité primitive récursive~» dépend du modèle de
calcul et des conventions de codage.  La formulation et la preuve
devront être vérifiées lors de la rédaction de la partie consacrée aux
machines de Turing.}

En particulier, on a l'inclusion
\[
\mathbf{EXPTIME}\subseteq\mathbf{PR}.
\]

\thingy La fonction d'Ackermann montre néanmoins que toutes les
fonctions calculables totales ne sont pas primitives récursives.  Pour
$m=2$, elle vérifie
\[
\begin{aligned}
A(2,n,0)&=2+n,\\
A(2,1,k+1)&=2,\\
A(2,n+1,k+1)&=A(2,\,A(2,n,k+1),\,k).
\end{aligned}
\]
Cette définition n'est pas une récursion primitive~: dans la dernière
ligne, la valeur de $k$ diminue lors de l'appel extérieur, mais un
appel imbriqué à $A$ intervient dans le deuxième argument.  Il ne
s'agit donc pas d'une simple récursion primitive sur une variable.

On peut montrer que, pour toute fonction primitive récursive
$f\colon\mathbb{N}^k\to\mathbb{N}$, il existe un entier $r=r(f)$ tel
que
\[
f(x_1,\ldots,x_k)
 \leq A(2,x_1+\cdots+x_k+3,r).
\]
Par conséquent, la fonction diagonale
\[
r\mapsto A(2,r,r)
\]
n'est pas primitive récursive.  Elle est pourtant définie par un
algorithme clair qui termine pour toute entrée.

\review{La démonstration de la majoration des fonctions primitives
récursives par les niveaux de la fonction d'Ackermann n'est pas donnée
dans les diapositives.  Elle est seulement annoncée et n'est donc pas
développée ici.}

\subsection{Numérotation des fonctions primitives récursives}
\label{subsection-numbering-primitive-recursive-functions}

\thingy Nous voulons maintenant coder les programmes primitifs
récursifs par des entiers.  Pour certains entiers $e\in\mathbb{N}$,
nous définirons une fonction
\[
\psi_e^{(k)}\colon\mathbb{N}^k\dasharrow\mathbb{N},
\]
appelée la fonction d'arité $k$ \defin[codée
(fonction)]{codée} par $e$.  L'entier $e$ doit être compris comme un
code source ou un programme.

Toute fonction primitive récursive
$f\colon\mathbb{N}^k\to\mathbb{N}$ sera égale à
$\psi_e^{(k)}$ pour au moins un entier $e$.  Cet entier décrit la
manière dont $f$ est construite à l'aide des clauses de la
définition~\ref{primitive-recursive-definition}.  Il n'est pas unique~:
une même fonction peut être calculée par de nombreux programmes
distincts,
\[
f=\psi_{e_1}^{(k)}
 =\psi_{e_2}^{(k)}
 =\cdots.
\]
L'indice $e$ décrit l'intention, tandis que la fonction $f$ est
l'extension.

\thingy Avec le codage de listes de~\ref{godel-coding}, on définit
$\psi_e^{(k)}$ par induction sur la construction primitive récursive.
\begin{itemize}
\item Si $e=\dbllangle0,k,i\dblrangle$, alors
\[
\psi_e^{(k)}(x_1,\ldots,x_k)=x_i.
\]
\item Si $e=\dbllangle1,k,c\dblrangle$, alors
\[
\psi_e^{(k)}(x_1,\ldots,x_k)=c.
\]
\item Si $e=\dbllangle2\dblrangle$, alors
\[
\psi_e^{(1)}(x)=x+1.
\]
\item Si
\[
e=\dbllangle3,k,d,c_1,\ldots,c_\ell\dblrangle,
\qquad
g_i:=\psi_{c_i}^{(k)},\qquad h:=\psi_d^{(\ell)},
\]
alors
\[
\psi_e^{(k)}(\underline{x})
 =h(g_1(\underline{x}),\ldots,g_\ell(\underline{x})).
\]
\item Si $e=\dbllangle4,k,d,c\dblrangle$, avec
$g:=\psi_c^{(k)}$ et $h:=\psi_d^{(k+2)}$, alors
\[
\begin{aligned}
\psi_e^{(k+1)}(\underline{x},0)
 &=g(\underline{x}),\\
\psi_e^{(k+1)}(\underline{x},z+1)
 &=h(\underline{x},
     \psi_e^{(k+1)}(\underline{x},z),z).
\end{aligned}
\]
\end{itemize}
Dans tous les autres cas, $\psi_e^{(k)}$ n'est pas définie.

Par définition, une fonction
$f\colon\mathbb{N}^k\dasharrow\mathbb{N}$ est primitive récursive si
et seulement s'il existe $e\in\mathbb{N}$ tel que
\[
f=\psi_e^{(k)}.
\]

\thingy La numérotation précédente permet de manipuler
algorithmiquement les programmes.  L'opération fondamentale suivante
consiste à fixer certains de leurs arguments.

\begin{thm}[Théorème s-m-n, version primitive récursive]
\label{smn-primitive-recursive}
Pour tous $m,n\in\mathbb{N}$, il existe une fonction primitive
récursive
\[
s_{m,n}\colon\mathbb{N}^{m+1}\to\mathbb{N}
\]
telle que
\[
\psi^{(n)}_{s_{m,n}(e,x_1,\ldots,x_m)}
       (y_1,\ldots,y_n)
 =
\psi^{(m+n)}_e
       (x_1,\ldots,x_m,y_1,\ldots,y_n)
\]
pour tous $e,\underline{x},\underline{y}$.
\end{thm}

\begin{proof}
Avec les conventions de codage précédentes, on peut prendre
\[
\begin{aligned}
s_{m,n}(e,\underline{x})
 :=\dbllangle
 &3,n,e,
 \dbllangle1,n,x_1\dblrangle,\ldots,
 \dbllangle1,n,x_m\dblrangle,\\
 &\dbllangle0,n,1\dblrangle,\ldots,
 \dbllangle0,n,n\dblrangle
 \dblrangle.
\end{aligned}
\]
Ce code décrit la composition du programme $e$ avec, pour ses $m$
premiers arguments, les fonctions constantes de valeurs
$x_1,\ldots,x_m$, et, pour ses $n$ derniers arguments, les projections.
Toutes les opérations de codage utilisées sont primitives récursives,
donc $s_{m,n}$ l'est aussi.
\end{proof}

En clair, $s_{m,n}$ transforme un programme prenant $m+n$ arguments en
un programme dans lequel les $m$ premiers arguments ont été fixés aux
valeurs $x_1,\ldots,x_m$, tandis que les $n$ derniers restent variables.

\subsection{Auto-référence et théorème de récursion}
\label{subsection-kleene-recursion-theorem-pr}

\thingy L'\defin[Quine (astuce de)]{astuce de Quine} est une technique
permettant à un programme d'avoir accès à son propre code source.  Son
idée peut être illustrée par la phrase autoréférente suivante~:

\begin{center}
\emph{Les mots suivants suivis des mêmes mots entre guillemets forment
une phrase intéressante~: «~les mots suivants suivis des mêmes mots
entre guillemets forment une phrase intéressante~».}
\end{center}

La première partie de la phrase sert de modèle et la seconde lui
fournit une représentation d'elle-même.  Une construction analogue
permet d'écrire un programme qui transmet son propre code source à une
fonction arbitraire.  Ainsi, même si un langage ne fournit aucune
instruction primitive donnant accès au code du programme courant, il
est possible de simuler cet accès.

\review{Le pseudocode de la diapositive est volontairement informel.
Il pourra être réintroduit sous forme d'exercice lors de l'intégration
des feuilles d'exercices.}

\begin{thm}[Kleene, version primitive récursive]
\label{kleene-recursion-theorem-pr}
Si $h\colon\mathbb{N}^{k+1}\dasharrow\mathbb{N}$ est primitive
récursive, il existe un entier $e$ tel que
\[
(\forall\underline{x})\qquad
\psi_e^{(k)}(\underline{x})=h(e,\underline{x}).
\]
Plus précisément, il existe une fonction primitive récursive
$b\colon\mathbb{N}^2\to\mathbb{N}$ telle que, si $e:=b(k,d)$, alors
\[
(\forall d)(\forall\underline{x})\qquad
\psi_e^{(k)}(\underline{x})
 =\psi_d^{(k+1)}(e,\underline{x}).
\]
\end{thm}

\begin{proof}
Soit $s:=s_{1,k}$ la fonction fournie par le théorème s-m-n.  La
fonction
\[
(t,\underline{x})\mapsto h(s(t,t),\underline{x})
\]
est primitive récursive.  Il existe donc un entier $c$ tel que
\[
\psi_c^{(k+1)}(t,\underline{x})
 =h(s(t,t),\underline{x}).
\]
Posons $e:=s(c,c)$.  Le théorème s-m-n donne alors
\[
\begin{aligned}
\psi_e^{(k)}(\underline{x})
 &=\psi_{s(c,c)}^{(k)}(\underline{x})\\
 &=\psi_c^{(k+1)}(c,\underline{x})\\
 &=h(s(c,c),\underline{x})\\
 &=h(e,\underline{x}).
\end{aligned}
\]
Les constructions qui, à $d$, associent un indice $c$, puis à $c$
l'indice $s(c,c)$, peuvent être effectuées primitivement récursivement,
ce qui donne la fonction uniforme $b$ annoncée.
\end{proof}

La moralité est que l'on peut donner à un programme accès à son propre
numéro, donc à une représentation de son code source, sans augmenter
la classe des fonctions qu'il peut calculer.

\thingy Il n'existe cependant pas d'interpréteur universel
\emph{primitif récursif} pour les programmes primitifs récursifs.

\begin{thm}\label{primitive-recursive-no-universality}
Il n'existe pas de fonction primitive récursive
$u\colon\mathbb{N}^2\to\mathbb{N}$ telle que
\[
u(e,x)=\psi_e^{(1)}(x)
\]
chaque fois que $\psi_e^{(1)}(x)\downarrow$.
\end{thm}

\begin{proof}
Supposons par l'absurde qu'une telle fonction $u$ existe.  La fonction
\[
(e,x)\mapsto u(e,x)+1
\]
est alors primitive récursive.  Par le théorème de récursion de Kleene,
il existe un indice $e$ tel que
\[
\psi_e^{(1)}(x)=u(e,x)+1
\]
pour tout $x$.  Mais la propriété supposée de $u$ donne simultanément
\[
u(e,x)=\psi_e^{(1)}(x),
\]
d'où
\[
u(e,x)=u(e,x)+1,
\]
ce qui est impossible.
\end{proof}

Autrement dit, un interpréteur du langage primitif récursif ne peut pas
être lui-même primitif récursif.  C'est un argument diagonal, analogue
à celui de Cantor.  Il repose à la fois sur le théorème s-m-n et sur le
fait que les fonctions primitives récursives sont totales.

On peut également retrouver cette absence d'universalité à l'aide de
la fonction d'Ackermann.  Pour chaque $k$ fixé, la fonction
\[
(m,n)\mapsto A(m,n,k)
\]
est primitive récursive, et on peut calculer primitivement
récursivement, en fonction de $k$, un indice $a(k)$ tel que
\[
\psi_{a(k)}^{(2)}(m,n)=A(m,n,k).
\]
Si la fonction universelle $(e,n)\mapsto\psi_e^{(1)}(n)$ était
primitive récursive, le théorème s-m-n permettrait d'en déduire que
$(n,k)\mapsto A(2,n,k)$ est primitive récursive, ce qui contredit les
limitations établies précédemment.

\subsection{Fonctions générales récursives}
\label{subsection-general-recursive-functions}

\thingy Les fonctions primitives récursives ne couvrent pas tous les
algorithmes.  Les programmes primitifs récursifs terminent toujours,
n'autorisent que des boucles bornées à l'avance et ne peuvent pas
s'interpréter eux-mêmes.  Pour obtenir une classe correspondant à la
calculabilité générale, il faut autoriser des recherches non bornées,
donc accepter que certains calculs ne terminent pas.

\begin{defn}\label{definition-mu-operator}
Soit
\[
g\colon\mathbb{N}^{k+1}\dasharrow\mathbb{N}.
\]
La fonction $\mu g\colon\mathbb{N}^k\dasharrow\mathbb{N}$ est définie
en posant que $\mu g(\underline{x})$ est le plus petit entier $z$ tel
que
\[
g(z,\underline{x})=0
\]
et tel que $g(i,\underline{x})\downarrow$ pour tout $0\leq i<z$, s'il
existe un tel entier.

Autrement dit, $\mu g(\underline{x})=z$ signifie
\[
g(z,\underline{x})=0
\quad\text{et}\quad
g(i,\underline{x})>0
\quad\text{pour tout $0\leq i<z$}.
\]
\end{defn}

Informatiquement, $\mu g$ correspond à la recherche non bornée
\[
\texttt{i=0; while (true) \{ if (g(i,x)==0) return i; i++; \}}.
\]
Si aucune valeur convenable n'est trouvée, ou si l'évaluation de
$g(i,\underline{x})$ ne termine pas avant d'atteindre une éventuelle
solution, le calcul ne termine pas.

\begin{defn}\label{general-recursive-definition}
La classe $\mathbf{R}$ des \defin[générale récursive
(fonction)]{fonctions générales récursives}, ou simplement
\defin[récursive (fonction)]{fonctions récursives}, est la plus petite
classe de fonctions
$\mathbb{N}^k\dasharrow\mathbb{N}$, pour $k$ variable, qui~:
\begin{itemize}
\item contient les projections, les constantes et le successeur~;
\item est stable par composition~;
\item est stable par récursion primitive~;
\item est stable par l'opérateur $\mu$.
\end{itemize}
\end{defn}

Cette dernière clause autorise précisément les boucles dont le nombre
d'itérations n'est pas borné à l'avance.  Les fonctions ainsi obtenues
peuvent être partielles.

\thingy On numérote les fonctions générales récursives exactement
comme les fonctions primitives récursives, mais en ajoutant une clause
pour l'opérateur $\mu$.  Pour certains entiers $e$, on définit donc
\[
\varphi_e^{(k)}\colon\mathbb{N}^k\dasharrow\mathbb{N}.
\]
Les codes correspondant aux projections, constantes, successeur,
composition et récursion primitive sont les mêmes que précédemment.
On ajoute la clause suivante~: si
\[
e=\dbllangle5,k,c\dblrangle
\quad\text{et}\quad
g:=\varphi_c^{(k+1)},
\]
alors
\[
\varphi_e^{(k)}=\mu g.
\]
Les autres codes ne définissent aucune fonction.

Une fonction $f\colon\mathbb{N}^k\dasharrow\mathbb{N}$ est récursive
si et seulement s'il existe un indice $e$ tel que
\[
f=\varphi_e^{(k)}.
\]

\thingy Contrairement au cas primitif récursif, il existe cette fois
une fonction récursive universelle~:
\[
(e,\underline{x})\mapsto\varphi_e^{(k)}(\underline{x}).
\]
Autrement dit, il existe un indice $u_k$ tel que
\[
\varphi_{u_k}^{(k+1)}(e,\underline{x})
 =\varphi_e^{(k)}(\underline{x}).
\]
Informatiquement, cette fonction est un interpréteur du langage des
fonctions récursives écrit dans ce même langage.

Pour justifier cette propriété, on introduit des \defin{arbres de
calcul}.  Un arbre de calcul de
$\varphi_e^{(k)}(\underline{x})$ ayant pour résultat $y$ est un arbre
fini dont la racine affirme
\[
\varphi_e^{(k)}(\underline{x})=y
\]
et dont les sous-arbres justifient les calculs intermédiaires requis
par la clause de construction correspondant à $e$.

Pour une composition, les sous-arbres établissent les valeurs des
fonctions intérieures puis celle de la fonction extérieure.  Pour une
récursion primitive, ils établissent successivement les valeurs
précédentes.  Pour une minimisation $\mu g$, ils établissent les
valeurs positives
\[
g(0,\underline{x}),\ldots,g(y-1,\underline{x})
\]
et la valeur nulle $g(y,\underline{x})$.

On a alors
\[
\varphi_e^{(k)}(\underline{x})=y
\]
si et seulement s'il existe un arbre de calcul fini qui l'atteste.
Or vérifier qu'un entier code un arbre de calcul valable est une
opération primitive récursive.  Il suffit donc de parcourir les
entiers $n=0,1,2,\ldots$, de tester si $n$ code un arbre de calcul
valable pour $\varphi_e^{(k)}(\underline{x})$, et, dès que c'est le
cas, de renvoyer la valeur inscrite à la racine.  Cette recherche non
bornée utilise une seule application de l'opérateur $\mu$ et définit
la fonction universelle.

\begin{thm}[Forme normale de Kleene]
\label{normal-form-theorem}
Il existe un prédicat primitif récursif $T$ sur $\mathbb{N}^3$ et une
fonction primitive récursive $U\colon\mathbb{N}\to\mathbb{N}$ tels que
\[
\varphi_e^{(k)}(x_1,\ldots,x_k)
 =
U\bigl(\mu T(e,\dbllangle x_1,\ldots,x_k\dblrangle)\bigr).
\]
\end{thm}

Ici, $T(n,e,\dbllangle x_1,\ldots,x_k\dblrangle)$ exprime que $n$
code un arbre de calcul valable de
$\varphi_e^{(k)}(x_1,\ldots,x_k)$, tandis que $U$ extrait de cet arbre
son résultat.  Ce théorème affirme donc que toute fonction récursive
peut être calculée à l'aide d'une unique recherche non bornée, entourée
de calculs primitifs récursifs.

\thingy Le théorème s-m-n reste valable pour la numérotation
$\varphi_e^{(k)}$.

\begin{thm}[Théorème s-m-n]\label{smn-recursive}
Pour tous $m,n$, il existe une fonction primitive récursive
\[
s_{m,n}\colon\mathbb{N}^{m+1}\to\mathbb{N}
\]
telle que
\[
\varphi^{(n)}_{s_{m,n}(e,x_1,\ldots,x_m)}
       (y_1,\ldots,y_n)
 =
\varphi^{(m+n)}_e
       (x_1,\ldots,x_m,y_1,\ldots,y_n).
\]
\end{thm}

Le fait que les fonctions manipulées puissent maintenant être
partielles ne change pas la construction syntaxique du programme.
La fonction $s_{m,n}$ reste donc primitive récursive.

\thingy Le théorème de récursion de Kleene s'étend lui aussi sans
modification essentielle.

\begin{thm}[Kleene]\label{kleene-recursion-theorem}
Si
\[
h\colon\mathbb{N}^{k+1}\dasharrow\mathbb{N}
\]
est récursive, il existe un indice $e$ tel que
\[
(\forall\underline{x})\qquad
\varphi_e^{(k)}(\underline{x})=h(e,\underline{x}).
\]
De plus, on peut calculer primitivement récursivement un tel indice à
partir d'un indice de $h$ et de l'arité $k$.
\end{thm}

\begin{proof}
La preuve est la même que dans le cas primitif récursif.  Soit
$s:=s_{1,k}$.  La fonction
\[
(t,\underline{x})\mapsto h(s(t,t),\underline{x})
\]
est récursive~; écrivons-la sous la forme
$\varphi_c^{(k+1)}(t,\underline{x})$.  Alors
\[
\begin{aligned}
\varphi_{s(c,c)}^{(k)}(\underline{x})
 &=\varphi_c^{(k+1)}(c,\underline{x})\\
 &=h(s(c,c),\underline{x}).
\end{aligned}
\]
L'indice $e:=s(c,c)$ convient.
\end{proof}

Une reformulation utile fait intervenir l'universalité.

\begin{thm}[Kleene--Rogers]\label{kleene-rogers-fixed-point}
Si $F\colon\mathbb{N}\to\mathbb{N}$ est récursive et si
$k\in\mathbb{N}$, il existe un indice $e$ tel que
\[
\varphi_e^{(k)}=\varphi_{F(e)}^{(k)}.
\]
\end{thm}

Ainsi, pour toute transformation calculable $F$ appliquée au code
source d'un programme, il existe un programme qui calcule la même
fonction que son transformé par $F$.

\thingy Malgré leur nom, les fonctions générales récursives ne sont pas
définies directement par des appels récursifs.  Le langage de la
définition~\ref{general-recursive-definition} ne contient que des
opérations élémentaires, la composition, la récursion primitive et la
recherche non bornée.  Le théorème de récursion permet néanmoins
d'implémenter les appels récursifs~: on donne à la fonction en cours de
définition accès à son propre indice, puis chaque appel récursif est
remplacé par un appel à la fonction universelle sur cet indice.

\subsection{Le problème de l'arrêt}
\label{subsection-halting-problem-recursive-functions}

\thingy Étant donné un programme d'indice $e$ et une entrée $x$, peut-on
déterminer algorithmiquement si le calcul
$\varphi_e^{(1)}(x)$ termine~?  La réponse est négative.

\begin{thm}[Turing]\label{undecidability-halting-problem}
Il n'existe pas de fonction récursive totale
$h\colon\mathbb{N}^2\to\mathbb{N}$ telle que
\[
h(e,x)=
\begin{cases}
1&\text{si $\varphi_e^{(1)}(x)\downarrow$,}\\
0&\text{si $\varphi_e^{(1)}(x)\uparrow$.}
\end{cases}
\]
\end{thm}

\begin{proof}
Supposons par l'absurde qu'une telle fonction $h$ existe.  Définissons
la fonction partielle
\[
v(e,x):=
\begin{cases}
42&\text{si $h(e,x)=0$,}\\
\uparrow&\text{si $h(e,x)=1$.}
\end{cases}
\]
Cette fonction est récursive~: on calcule $h(e,x)$, puis on renvoie
$42$ dans le premier cas et on effectue une recherche infinie dans le
second.

Par le théorème de récursion de Kleene, il existe un indice $e$ tel que
\[
\varphi_e^{(1)}(x)=v(e,x).
\]
Si $\varphi_e^{(1)}(x)\downarrow$, alors $h(e,x)=1$, donc
$v(e,x)\uparrow$, ce qui contredit l'égalité précédente.  Si
$\varphi_e^{(1)}(x)\uparrow$, alors $h(e,x)=0$, donc
$v(e,x)\downarrow$, ce qui donne encore une contradiction.
\end{proof}

Une version plus directement diagonale consiste à supposer que $h$
décide l'arrêt et à définir
\[
v(e):=
\begin{cases}
42&\text{si $h(e,e)=0$,}\\
\uparrow&\text{si $h(e,e)=1$.}
\end{cases}
\]
Si $v=\varphi_c^{(1)}$, alors l'étude de $\varphi_c^{(1)}(c)$ conduit
immédiatement à une contradiction~: elle termine si et seulement si
elle ne termine pas.

\thingy Les fonctions primitives récursives sont toutes totales, de
sorte que leur problème de l'arrêt est trivial.  Elles ne possèdent en
revanche pas de fonction universelle primitive récursive.  Les
fonctions générales récursives peuvent être partielles et disposent
d'une fonction universelle, mais leur problème de l'arrêt est
indécidable.  L'universalité est précisément ce qui permet de
construire l'argument diagonal.

\subsection{Machines de Turing}
\label{subsection-turing-machines}

\thingy Une \defin[Turing (machine de)]{machine de Turing} est une
modélisation d'un ordinateur extrêmement simple.  Elle possède un état
interne choisi dans un ensemble fini et une mémoire illimitée sous la
forme d'une bande linéaire divisée en cellules.  Chaque cellule contient
un symbole pris dans un alphabet fini.  Une tête de lecture et
d'écriture se trouve au-dessus d'une unique cellule à chaque instant.

Le programme de la machine indique, en fonction de l'état courant et
du symbole lu par la tête~:
\begin{itemize}
\item le nouvel état~;
\item le symbole à écrire dans la cellule courante~;
\item la direction dans laquelle déplacer la tête, à gauche ou à
  droite.
\end{itemize}
La machine poursuit son exécution jusqu'à atteindre un état spécial
d'arrêt.

\begin{defn}\label{definition-turing-machine}
Une machine de Turing déterministe à une bande, deux symboles et
$m\geq2$ états est la donnée~:
\begin{itemize}
\item d'un ensemble fini $Q$ de cardinal $m$, identifié à
  $\{0,\ldots,m-1\}$~;
\item d'un ensemble de symboles de bande
  $\Sigma=\{0,1\}$~;
\item d'une fonction
\[
\delta\colon
(Q\setminus\{0\})\times\Sigma
 \longrightarrow
Q\times\Sigma\times\{\texttt{L},\texttt{R}\},
\]
appelée le programme de la machine.
\end{itemize}
L'état $0$ est l'état d'arrêt.
\end{defn}

Une \defin{configuration} d'une telle machine est un triplet
$(q,\beta,i)$, où~:
\begin{itemize}
\item $q\in Q$ est l'état courant~;
\item $\beta\colon\mathbb{Z}\to\Sigma$ est la bande, supposée ne prendre
  qu'un nombre fini de valeurs différentes de $0$~;
\item $i\in\mathbb{Z}$ est la position de la tête.
\end{itemize}

Si $q\neq0$ et si
\[
(q',y,d)=\delta(q,\beta(i)),
\]
la configuration suivante est $(q',\beta',i')$, où
\[
i'=
\begin{cases}
i-1&\text{si $d=\texttt{L}$,}\\
i+1&\text{si $d=\texttt{R}$,}
\end{cases}
\]
et
\[
\beta'(j)=
\begin{cases}
y&\text{si $j=i$,}\\
\beta(j)&\text{si $j\neq i$.}
\end{cases}
\]

\thingy À partir d'une configuration initiale $C$, la
\defin[exécution (trace d')]{trace d'exécution} est la suite finie ou
infinie
\[
C^{(0)},C^{(1)},C^{(2)},\ldots
\]
où $C^{(0)}=C$ et où $C^{(n+1)}$ est la configuration suivant
$C^{(n)}$, tant que l'état courant n'est pas $0$.  Si une configuration
$C^{(n)}$ d'état $0$ est atteinte, la machine s'arrête après $n$
étapes et $C^{(n)}$ est sa configuration finale.

\thingy Une machine de Turing peut être codée par un entier, puisque
son programme est une table finie.  Une configuration peut elle aussi
être codée par un entier, puisque la bande ne contient qu'un nombre
fini de symboles non nuls.

La fonction
\[
(M,C)\mapsto C'
\]
associant à une machine et à une configuration la configuration
suivante est primitive récursive.  Par récursion primitive, la
fonction
\[
(n,M,C)\mapsto C^{(n)}
\]
associant la configuration atteinte après $n$ étapes l'est également.

En revanche, la fonction qui associe à $(M,C)$ sa configuration finale
si la machine s'arrête, et qui est non définie sinon, est générale
récursive.  En effet, on peut rechercher sans borne le premier entier
$n$ tel que $C^{(n)}$ soit une configuration d'arrêt.

\thingy Pour définir précisément les fonctions calculées par les
machines de Turing, fixons une convention d'entrée et de sortie.  Une
fonction
\[
f\colon\mathbb{N}^k\dasharrow\mathbb{N}
\]
est dite \defin[calculable par machine de Turing]{calculable par
machine de Turing} lorsqu'il existe une machine qui, sur les entrées
$x_1,\ldots,x_k$, part d'une configuration initiale dans laquelle~:
\begin{itemize}
\item l'état vaut $1$ et la tête se trouve en position $0$~;
\item la partie négative du ruban est laissée intacte~;
\item la partie positive contient le mot
\[
0\,1^{x_1}\,0\,1^{x_2}\,0\cdots0\,1^{x_k}\,0,
\]
suivi d'une infinité de $0$~;
\end{itemize}
et telle que~:
\begin{itemize}
\item si $f(x_1,\ldots,x_k)\uparrow$, la machine ne s'arrête pas~;
\item si $f(x_1,\ldots,x_k)\downarrow=y$, elle s'arrête avec la tête en
  position $0$, sans modifier la partie négative du ruban, et avec
  \[
  0\,1^y\,0
  \]
  sur la partie positive du ruban, suivi d'une infinité de $0$.
\end{itemize}

Ces conventions sont arbitraires mais sans importance pour la classe
de fonctions obtenue.

\begin{thm}\label{recursive-functions-turing-machines}
Une fonction
$f\colon\mathbb{N}^k\dasharrow\mathbb{N}$ est générale récursive si et
seulement si elle est calculable par une machine de Turing.
\end{thm}

\begin{proof}
Pour simuler les machines de Turing par les fonctions récursives, on
code les programmes et les configurations par des entiers.  La
fonction donnant la configuration après $n$ étapes est primitive
récursive.  Une recherche non bornée permet de trouver le premier
$n$ pour lequel la machine est à l'arrêt, puis une fonction primitive
récursive décode la sortie.

Réciproquement, on montre par induction sur la
définition~\ref{general-recursive-definition} que toute fonction
générale récursive est calculable par machine de Turing.  Il faut
programmer les projections, les constantes et le successeur, puis
montrer que les fonctions calculables par machine de Turing sont
stables par composition, récursion primitive et minimisation.  Ces
constructions reviennent à sauvegarder les entrées et les résultats
intermédiaires sur la bande, puis à exécuter successivement les
machines correspondant aux sous-calculs.
\end{proof}

\review{Les diapositives qualifient la seconde direction de
«~fastidieuse mais pas difficile~» et n'en donnent pas tous les détails.
La preuve ci-dessus reste donc une esquisse à développer ou à renvoyer
à un exercice.}

Cette équivalence est constructive~: on peut calculer un code de
machine de Turing à partir d'un indice de fonction récursive, et
réciproquement.

\thingy De nombreuses variantes de la définition sont possibles~:
machines à plusieurs bandes, plusieurs têtes, davantage de symboles ou
transitions non déterministes.  Du point de vue de la calculabilité,
ces variantes ne rendent pas le modèle plus puissant, sauf cas
dégénérés.  Elles peuvent cependant modifier considérablement la
complexité des simulations.

Une machine de Turing universelle peut recevoir sur sa bande le code
d'une autre machine $M$ et d'une configuration initiale $C$, puis
simuler l'exécution de $M$ sur $C$.  Elle s'arrête si et seulement si
$M$ s'arrête.  Cette propriété est la traduction, pour les machines de
Turing, de l'existence d'une fonction récursive universelle.

\begin{thm}\label{halting-problem-turing-machines}
Il n'existe pas d'algorithme qui, prenant en entrée une machine de
Turing $M$ et une configuration $C$, décide si $M$ s'arrête à partir
de $C$.
\end{thm}

\begin{proof}
Si un tel algorithme existait, on pourrait coder $M$ et $C$ par des
entiers et décider le problème de l'arrêt pour les fonctions
récursives, en contradiction avec le
théorème~\ref{undecidability-halting-problem}.
\end{proof}

Le problème reste indécidable si l'on demande seulement si une machine
s'arrête à partir d'une bande vierge.  En effet, étant donnés $M$ et
$C$, le théorème s-m-n permet de construire une machine $M^*$ qui
prépare d'abord la configuration $C$ à partir de la bande vierge, puis
simule $M$.  La machine $M^*$ s'arrête sur bande vierge si et seulement
si $M$ s'arrête sur $C$.

\thingy La \defin[castor affairé]{fonction du castor affairé} associe à
$m$ le nombre maximal $B(m)$ d'étapes effectuées, à partir d'une bande
vierge, par une machine de Turing ayant au plus $m$ états et qui finit
par s'arrêter.  Ce maximum existe, puisqu'il n'existe qu'un nombre fini
de machines ayant au plus $m$ états.

\begin{prop}\label{busy-beaver-not-computable}
La fonction $B$ n'est pas calculable.
\end{prop}

\begin{proof}
Supposons qu'il existe une fonction calculable
$f\colon\mathbb{N}\to\mathbb{N}$ telle que
\[
B(m)\leq f(m)
\]
pour tout $m$.  Étant donnée une machine $M$ à $m$ états, on pourrait
l'exécuter pendant $f(m)$ étapes.  Si elle s'arrête pendant ce temps,
on répondrait qu'elle termine.  Sinon, la définition de $B(m)$ assure
qu'elle ne terminera jamais.  On aurait ainsi un algorithme décidant
le problème de l'arrêt sur bande vierge, ce qui est impossible.
\end{proof}

La fonction du castor affairé croît donc plus vite que toute fonction
calculable, au sens suivant~: pour toute fonction calculable totale
$f\colon\mathbb{N}\to\mathbb{N}$, il existe $m_0$ tel que
\[
B(m)>f(m)
\]
pour tout $m\geq m_0$.

\review{Les diapositives esquissent une preuve de cette domination à
l'aide d'une famille de machines dont le nombre d'états est contrôlé.
La formulation uniforme et les détails de codage devront être
vérifiés.}

\subsection{Décidabilité et semi-décidabilité}
\label{subsection-decidability-semidecidability}

\thingy D'après les équivalences précédentes, une fonction partielle
est calculable si et seulement si elle est générale récursive, si et
seulement si elle est calculable par machine de Turing.  Nous
emploierons librement ces différentes descriptions.

On rencontre aussi les synonymes suivants~:
\begin{itemize}
\item une fonction partielle calculable est parfois dite
  «~semi-calculable~» ou «~générale récursive~»~;
\item un ensemble décidable est parfois dit «~calculable~» ou
  «~récursif~»~;
\item un ensemble semi-décidable est parfois dit
  «~semi-calculable~», «~calculablement énumérable~» ou
  «~récursivement énumérable~».
\end{itemize}

\begin{prop}\label{decidable-iff-semidecidable-and-complement}
Une partie $A\subseteq\mathbb{N}^k$ est décidable si et seulement si
$A$ et son complémentaire
\[
\complement A:=\mathbb{N}^k\setminus A
\]
sont tous deux semi-décidables.
\end{prop}

\begin{proof}
Si $A$ est décidable, son complémentaire est décidable en échangeant
les réponses $0$ et $1$.  De plus, tout ensemble décidable est
semi-décidable~: il suffit d'exécuter son algorithme de décision et de
boucler indéfiniment lorsque celui-ci répond «~non~».

Réciproquement, supposons que $A$ et $\complement A$ soient
semi-décidables.  Soient $e_1$ et $e_2$ des indices tels que
\[
\varphi_{e_1}(\underline{x})\downarrow
 \quad\Longleftrightarrow\quad
\underline{x}\in A
\]
et
\[
\varphi_{e_2}(\underline{x})\downarrow
 \quad\Longleftrightarrow\quad
\underline{x}\notin A.
\]
On exécute les deux calculs en parallèle, en alternant leurs étapes.
Exactement l'un des deux finit par terminer.  S'il s'agit du premier,
on répond «~oui~»~; s'il s'agit du second, on répond «~non~».
L'algorithme ainsi obtenu termine sur toute entrée et décide $A$.
\end{proof}

Le «~lancement en parallèle~» ne nécessite aucun véritable parallélisme.
On peut simuler successivement une étape du premier programme, puis une
étape du second, et recommencer jusqu'à ce que l'un d'eux s'arrête.

\thingy Le \defin[problème de l'arrêt]{problème de l'arrêt} est
l'ensemble
\[
\mathscr{H}
 :=\{(e,x)\in\mathbb{N}^2:
           \varphi_e^{(1)}(x)\downarrow\}.
\]
Il n'est pas décidable, d'après le
théorème~\ref{undecidability-halting-problem}.  Il est en revanche
semi-décidable~: sur l'entrée $(e,x)$, il suffit d'exécuter
$\varphi_e^{(1)}(x)$ et de répondre «~oui~» si le calcul termine.

La proposition~\ref{decidable-iff-semidecidable-and-complement} montre
donc que son complémentaire $\complement\mathscr{H}$ n'est pas
semi-décidable.  Il est possible de reconnaître en temps fini un
calcul qui termine, en observant effectivement sa terminaison, mais il
n'existe aucun procédé analogue reconnaissant tous les calculs qui ne
termineront jamais.

\thingy Si $A\subseteq\mathbb{N}$ est décidable et si
$f\colon\mathbb{N}\to\mathbb{N}$ est calculable totale, alors l'image
\[
f(A):=\{f(i):i\in A\}
\]
est semi-décidable.

En effet, pour semi-décider si $m\in f(A)$, on parcourt
$i=0,1,2,\ldots$.  Pour chaque $i$, on décide si $i\in A$ et, si c'est
le cas, on calcule $f(i)$.  Si $f(i)=m$, on termine en répondant
«~oui~»~; sinon, on poursuit la recherche.

Réciproquement, tout ensemble semi-décidable $B\subseteq\mathbb{N}$
peut être obtenu comme image d'un ensemble décidable par une fonction
totale calculable.

Une autre formulation, lorsque $B$ est non vide, est la suivante~:
$B$ est semi-décidable si et seulement s'il existe une fonction totale
calculable $f\colon\mathbb{N}\to\mathbb{N}$ telle que
\[
f(\mathbb{N})=B.
\]
C'est ce qui justifie la terminologie «~calculablement énumérable~».

\thingy Les ensembles décidables sont stables par réunions finies,
intersections finies et complémentaire.  Ils ne sont cependant pas
stables par projection.  Le problème de l'arrêt en donne déjà un
exemple, puisqu'il peut être présenté comme la projection d'un ensemble
décidable exprimant qu'un calcul s'arrête en au plus $n$ étapes.

Les ensembles semi-décidables sont stables par réunions finies,
intersections finies et projections.  Ils ne sont pas stables par
complémentaire, puisque $\mathscr{H}$ est semi-décidable tandis que
$\complement\mathscr{H}$ ne l'est pas.

\subsection{Le théorème de Rice}
\label{subsection-rice-theorem}

\thingy Soit $\mathbf{R}^{(1)}$ l'ensemble des fonctions partielles
calculables $\mathbb{N}\dasharrow\mathbb{N}$, et considérons
l'application
\[
\Phi\colon e\mapsto\varphi_e^{(1)}.
\]
Cette application associe à un programme, son intention, la fonction
qu'il calcule, son extension.

\begin{thm}[Rice]\label{rice-theorem}
Soit $F\subseteq\mathbf{R}^{(1)}$ un ensemble de fonctions partielles
calculables.  Si
\[
\Phi^{-1}(F)
 =\{e\in\mathbb{N}:\varphi_e^{(1)}\in F\}
\]
est décidable, alors
\[
F=\varnothing
\qquad\text{ou}\qquad
F=\mathbf{R}^{(1)}.
\]
\end{thm}

Autrement dit, aucune propriété non triviale de la fonction calculée
par un programme n'est décidable en examinant le programme.  Par
exemple, aucun des ensembles suivants n'est décidable~:
\[
\{e:\varphi_e(0)\downarrow\},
\]
\[
\{e:\varphi_e\text{ est totale}\},
\]
ou encore
\[
\{e:(\forall n)\,
       (\varphi_e(n)\downarrow\limp\varphi_e(n)=0)\}.
\]

\begin{proof}
Supposons par l'absurde que $\Phi^{-1}(F)$ soit décidable, avec
$F\neq\varnothing$ et $F\neq\mathbf{R}^{(1)}$.  Choisissons
$f\in F$ et $g\notin F$.  Définissons
\[
h(e,x):=
\begin{cases}
f(x)&\text{si $e\notin\Phi^{-1}(F)$,}\\
g(x)&\text{si $e\in\Phi^{-1}(F)$.}
\end{cases}
\]
La fonction $h$ est calculable, puisque l'appartenance de $e$ à
$\Phi^{-1}(F)$ est supposée décidable.

Par le théorème de récursion de Kleene, il existe un indice $e$ tel
que
\[
\varphi_e^{(1)}(x)=h(e,x).
\]
Si $e\in\Phi^{-1}(F)$, alors $h(e,x)=g(x)$ pour tout $x$, donc
$\varphi_e^{(1)}=g\notin F$, contradiction.  Si
$e\notin\Phi^{-1}(F)$, alors $\varphi_e^{(1)}=f\in F$, contradiction
également.
\end{proof}

\thingy Le théorème de Rice peut aussi être démontré par réduction au
problème de l'arrêt.  Supposons $F$ non trivial et, quitte à remplacer
$F$ par son complémentaire, supposons que la fonction nulle part
définie n'appartienne pas à $F$.  Choisissons $f\in F$ et un indice
$a$ de $f$.

Pour $(e,x)\in\mathbb{N}^2$, construisons un programme d'indice
$b(e,x)$ qui commence par simuler $\varphi_e(x)$.  Si cette simulation
termine, il calcule ensuite $f$ sur son argument.  Ainsi,
\[
\varphi_{b(e,x)}=
\begin{cases}
f&\text{si $\varphi_e(x)\downarrow$,}\\
{\uparrow}&\text{si $\varphi_e(x)\uparrow$,}
\end{cases}
\]
où ${\uparrow}$ désigne la fonction nulle part définie.  On obtient
donc
\[
b(e,x)\in\Phi^{-1}(F)
\quad\Longleftrightarrow\quad
(e,x)\in\mathscr{H}.
\]
Si $\Phi^{-1}(F)$ était décidable, le problème de l'arrêt le serait
aussi, ce qui est impossible.

\subsection{Réductions}
\label{subsection-reductions}

\thingy Pour montrer qu'un problème $A$ est indécidable, on cherche
souvent à transformer toute instance d'un problème déjà connu comme
indécidable, par exemple $\mathscr{H}$, en une instance de $A$.  Cette
transformation est appelée une \defin{réduction}.

\begin{defn}\label{definition-many-one-reduction}
Si $A,B\subseteq\mathbb{N}$, on écrit
\[
A\mathrel{\leq_{\mathrm m}}B
\]
lorsqu'il existe une fonction totale calculable
$\rho\colon\mathbb{N}\to\mathbb{N}$ telle que
\[
\rho(m)\in B
\quad\Longleftrightarrow\quad
m\in A.
\]
On parle de \defin[réduction many-to-one]{réduction
\textit{many-to-one}}.
\end{defn}

Ainsi, pour décider si $m\in A$, il suffirait de calculer
$\rho(m)$ puis de décider si $\rho(m)\in B$.  Par conséquent, si
$A\leq_{\mathrm m}B$ et si $B$ est décidable, alors $A$ est décidable.
La même implication vaut pour la semi-décidabilité.  En particulier,
si
\[
\mathscr{H}\leq_{\mathrm m}D,
\]
alors $D$ n'est pas décidable.

La relation $\leq_{\mathrm m}$ est réflexive et transitive.  La relation
\[
A\equiv_{\mathrm m}B
\quad\Longleftrightarrow\quad
A\leq_{\mathrm m}B
\ \text{et}\
B\leq_{\mathrm m}A
\]
est donc une relation d'équivalence.

\thingy Une notion plus générale autorise un algorithme à interroger
plusieurs fois un oracle décidant $B$ au cours de son exécution.

Informellement, on écrit
\[
A\leq_{\mathrm T}B
\]
lorsqu'il existe un algorithme qui décide $A$ et qui peut, au cours de
son calcul, demander si un entier $n$ appartient à $B$.  Cette notion
est appelée \defin[Turing (réduction de)]{réduction de Turing}.

La réduction de Turing est plus faible que la réduction many-to-one~:
une réduction many-to-one ne pose qu'une question à l'oracle, à la fin
du calcul.  Une réduction de Turing peut en poser un nombre arbitraire,
mais fini, et choisir chaque nouvelle question en fonction des réponses
précédentes.

\review{Les diapositives proposent plusieurs formalisations possibles
des réductions de Turing.  Leur définition entièrement formelle devra
être harmonisée avec la présentation retenue pour les oracles dans le
document final.}

Si $A\leq_{\mathrm T}B$ et si $B$ est décidable, alors $A$ est
décidable.  Il s'ensuit que si $\mathscr{H}\leq_{\mathrm T}D$, alors
$D$ est indécidable.

\subsection{Le $\lambda$-calcul non typé}
\label{subsection-untyped-lambda-calculus}

\thingy Le \defin[$\lambda$-calcul]{$\lambda$-calcul non typé}
manipule des expressions telles que
\[
\lambda x.\lambda y.\lambda z.((xz)(yz)),
\qquad
\lambda f.\lambda x.f(f(f(f(fx)))),
\]
ou
\[
(\lambda x.(xx))(\lambda x.(xx)).
\]
Ces expressions, appelées \defin[$\lambda$-terme]{$\lambda$-termes},
représentent intuitivement des fonctions prenant d'autres fonctions en
entrée et renvoyant des fonctions.

Les deux constructions fondamentales sont~:
\begin{itemize}
\item l'\defin[application ($\lambda$-calcul)]{application} $(PQ)$,
  qui consiste à appliquer la fonction représentée par $P$ à
  l'argument représenté par $Q$~;
\item l'\defin[abstraction ($\lambda$-calcul)]{abstraction}
  $\lambda v.E$, qui représente la fonction associant à un argument
  la valeur obtenue en remplaçant $v$ par cet argument dans $E$.
\end{itemize}

\begin{defn}\label{definition-lambda-terms}
Un terme du $\lambda$-calcul est défini inductivement de la manière
suivante~:
\begin{itemize}
\item toute variable est un terme~;
\item si $P$ et $Q$ sont des termes, alors $(PQ)$ est un terme~;
\item si $v$ est une variable et $E$ un terme, alors $\lambda v.E$ est
  un terme.
\end{itemize}
Dans $\lambda v.E$, les occurrences de $v$ dans $E$ auxquelles ce
$\lambda$ s'applique sont dites \defin[liée
(variable)]{liées}.  Les autres occurrences de variables sont dites
\defin[libre (variable)]{libres}.  Un terme sans variable libre est dit
\defin[clos ($\lambda$-terme)]{clos}.
\end{defn}

Par convention, l'application est parenthésée vers la gauche~:
\[
xyz:=((xy)z).
\]
On abrège
\[
\lambda u.\lambda v.E
\]
en $\lambda uv.E$.  Enfin, l'abstraction est moins prioritaire que
l'application~:
\[
\lambda x.xy
\]
signifie $\lambda x.(xy)$ et non $(\lambda x.x)y$.

\thingy Le nom d'une variable liée ne doit pas avoir d'importance.  On
appelle \defin[$\alpha$-conversion]{$\alpha$-conversion} le renommage
cohérent d'une variable liée.  Ainsi,
\[
\lambda x.x\equiv\lambda y.y.
\]
Il faut cependant éviter de capturer une variable libre.  Par exemple,
\[
\lambda y.xy\not\equiv\lambda x.xx.
\]
Dans une expression comportant plusieurs abstractions portant le même
nom, une occurrence est liée par le $\lambda$ le plus intérieur
approprié.

Une façon de supprimer complètement les noms de variables liées
consiste à utiliser les indices de De Bruijn~: chaque occurrence liée
est remplacée par la distance qui la sépare de l'abstraction qui la
lie.  Deux termes sont alors $\alpha$-équivalents si et seulement si
leurs écritures à l'aide de ces indices sont identiques.

\begin{defn}\label{definition-beta-reduction}
Un \defin{redex} est un terme de la forme
\[
(\lambda v.E)T.
\]
Son \defin[réduit ($\lambda$-calcul)]{réduit} est le terme
\[
E[v\backslash T]
\]
obtenu en remplaçant dans $E$ les occurrences libres de $v$ par $T$,
en renommant au besoin des variables liées pour éviter toute capture.

Le remplacement d'un redex par son réduit est appelé une
\defin[$\beta$-réduction]{$\beta$-réduction}.
\end{defn}

Par exemple,
\[
(\lambda x.xx)y\longrightarrow yy,
\]
tandis que
\[
(\lambda x.xx)(\lambda x.xx)
 \longrightarrow
(\lambda x.xx)(\lambda x.xx),
\]
de sorte que ce dernier terme se réduit indéfiniment.

On note $T\longrightarrow T'$ lorsque $T'$ est obtenu par une seule
$\beta$-réduction, et
\[
T\twoheadrightarrow T'
\]
lorsque $T'$ est obtenu par une suite finie, éventuellement vide, de
$\beta$-réductions.

Un terme ne contenant aucun redex est dit en
\defin[forme normale ($\lambda$-calcul)]{forme normale}.  Il est
\defin[faiblement normalisable]{$\beta$-faiblement normalisable}
lorsqu'au moins une suite de réductions mène à une forme normale, et
\defin[fortement normalisable]{$\beta$-fortement normalisable} lorsque
toute suite de réductions termine.

\begin{thm}[Church--Rosser]\label{church-rosser-theorem}
Si
\[
T\twoheadrightarrow T'_1
\quad\text{et}\quad
T\twoheadrightarrow T'_2,
\]
alors il existe $T''$ tel que
\[
T'_1\twoheadrightarrow T''
\quad\text{et}\quad
T'_2\twoheadrightarrow T''.
\]
\end{thm}

En particulier, si $T'_1$ et $T'_2$ sont en forme normale, ils sont
$\alpha$-équivalents.  La forme normale, lorsqu'elle existe, est donc
unique à $\alpha$-conversion près.

\thingy Pour éviter d'avoir à démontrer ici le théorème de
Church--Rosser, on choisit une stratégie déterministe de réduction.  Le
\defin[redex extérieur gauche]{redex extérieur gauche} est le redex
dont le symbole $\lambda$ est situé le plus à gauche dans l'écriture du
terme.  On note
\[
T\longrightarrow_{\mathsf{lft}}T'
\]
la réduction de ce redex, et
\[
T\twoheadrightarrow_{\mathsf{lft}}T'
\]
une suite de telles réductions.

\begin{thm}[Curry et al.]\label{leftmost-normalization}
Si un terme $T$ possède une forme normale $T'$, alors
\[
T\twoheadrightarrow_{\mathsf{lft}}T'.
\]
\end{thm}

La stratégie extérieure gauche est donc déterministe et trouve une
forme normale chaque fois qu'une forme normale peut être atteinte.

\subsection{Entiers de Church et fonctions représentables}
\label{subsection-church-numerals}

\thingy Pour chaque entier $n\in\mathbb{N}$, on définit l'\defin[Church
(entier de)]{entier de Church}
\[
\overline{n}:=\lambda fx.f^{\circ n}(x).
\]
Ainsi,
\[
\begin{aligned}
\overline{0}&:=\lambda fx.x,\\
\overline{1}&:=\lambda fx.fx,\\
\overline{2}&:=\lambda fx.f(fx),\\
\overline{3}&:=\lambda fx.f(f(fx)).
\end{aligned}
\]
Intuitivement, $\overline{n}$ reçoit une fonction $f$ et un argument
$x$, puis applique $f$ exactement $n$ fois à $x$.

Posons
\[
A:=\lambda mfx.f(mfx).
\]
Alors
\[
A\overline{n}
 \twoheadrightarrow_{\mathsf{lft}}
\overline{n+1},
\]
de sorte que $A$ représente la fonction successeur.

\begin{defn}\label{definition-lambda-representable}
Une fonction
\[
f\colon\mathbb{N}^k\dasharrow\mathbb{N}
\]
est \defin[représentable dans le $\lambda$-calcul]{représentable par un
$\lambda$-terme} lorsqu'il existe un terme clos $t$ tel que, pour tous
$x_1,\ldots,x_k\in\mathbb{N}$~:
\begin{itemize}
\item si $f(x_1,\ldots,x_k)\downarrow=y$, alors
\[
t\overline{x_1}\cdots\overline{x_k}
 \twoheadrightarrow_{\mathsf{lft}}\overline{y};
\]
\item si $f(x_1,\ldots,x_k)\uparrow$, alors la réduction extérieure
  gauche de
  $t\overline{x_1}\cdots\overline{x_k}$ ne termine pas.
\end{itemize}
\end{defn}

À titre d'exemple,
\[
\lambda mfx.f(mfx)
\]
représente le successeur, tandis que
\[
\lambda mnfx.nf(mfx)
\]
représente l'addition, et
\[
\lambda mnf.n(mf)
\]
représente la multiplication.

On peut également représenter la fonction
\[
(m,n,p)\mapsto
\begin{cases}
m&\text{si $p=0$,}\\
n&\text{si $p\geq1$.}
\end{cases}
\]

\thingy Les projections, les constantes, le successeur et la
composition sont aisément représentables dans le $\lambda$-calcul.  La
représentation de la récursion primitive exige de coder des couples
d'entiers.

On peut poser
\[
\begin{aligned}
\overline{m,n}
 &:=\lambda fgx.f^{\circ m}(g^{\circ n}(x)),\\
\Pi&:=\lambda mnfgx.(mf)(ngx),\\
\pi_1&:=\lambda pfx.pf(\lambda z.z)x,\\
\pi_2&:=\lambda pgx.p(\lambda z.z)gx.
\end{aligned}
\]
Alors
\[
\Pi\overline{m}\,\overline{n}
 \twoheadrightarrow_{\mathsf{lft}}\overline{m,n},
\]
et
\[
\pi_1\overline{m,n}
 \twoheadrightarrow_{\mathsf{lft}}\overline{m},
\qquad
\pi_2\overline{m,n}
 \twoheadrightarrow_{\mathsf{lft}}\overline{n}.
\]

À l'aide de ces couples, on peut faire évoluer simultanément la valeur
d'une récursion primitive et son compteur.  Ceci permet de représenter
toute fonction primitive récursive par un $\lambda$-terme.

\review{La formule explicite donnée dans les diapositives pour
représenter la récursion primitive est longue et dépend des conventions
précises de codage des couples.  Elle devra être vérifiée caractère par
caractère avant d'être intégrée dans une version définitive.}

\thingy Pour représenter toutes les fonctions générales récursives, il
reste à implémenter des appels récursifs, ou de façon équivalente
l'opérateur $\mu$.  On utilise pour cela le
\defin[Curry (combinateur $\mathsf{Y}$ de)]{combinateur de point fixe
de Curry}
\[
\mathsf{Y}
 :=\lambda f.((\lambda x.f(xx))(\lambda x.f(xx))).
\]
En effet,
\[
\begin{aligned}
\mathsf{Y}f
 &\longrightarrow
 f((\lambda x.f(xx))(\lambda x.f(xx)))\\
 &=f(\mathsf{Y}f),
\end{aligned}
\]
à réduction près.  Le terme $\mathsf{Y}f$ se comporte donc comme un
point fixe de $f$.

Si l'on veut définir une fonction récursive par
\[
\texttt{let rec }h=E,
\]
où $E$ contient des occurrences de $h$, on remplace cette définition
par un terme construit à partir de
\[
\mathsf{Y}(\lambda h.E).
\]
Le combinateur $\mathsf{Y}$ fournit ainsi au corps $E$ une
représentation de la fonction en cours de définition.

Le fonctionnement de $\mathsf{Y}$ dépend toutefois de la stratégie de
réduction.  La variante
\[
\mathsf{Z}
 :=\lambda f.
 ((\lambda x.f(\lambda v.xxv))
  (\lambda x.f(\lambda v.xxv)))
\]
maintient le point fixe sous une abstraction et convient mieux aux
langages qui évaluent les arguments avant l'appel d'une fonction.

\thingy L'opérateur $\mu$ peut être représenté en programmant la
recherche récursive
\[
h(z,x_1,\ldots,x_k)=
\begin{cases}
z&\text{si $g(z,x_1,\ldots,x_k)=0$,}\\
h(z+1,x_1,\ldots,x_k)
 &\text{sinon.}
\end{cases}
\]
Alors
\[
\mu g(x_1,\ldots,x_k)=h(0,x_1,\ldots,x_k).
\]
Le combinateur de point fixe permet de représenter $h$, et donc
$\mu g$, dans le $\lambda$-calcul.

Réciproquement, un $\lambda$-terme peut être codé par un entier.  Le
test déterminant si un terme est en forme normale, ainsi que la
fonction calculant sa réduction extérieure gauche, sont primitifs
récursifs.  Par récursion primitive, on peut calculer le terme obtenu
après $n$ réductions.  Une recherche non bornée permet alors de trouver
la première forme normale, lorsqu'elle existe.  Les fonctions
récursives peuvent donc simuler le $\lambda$-calcul.

\begin{thm}\label{recursive-functions-lambda-calculus}
Une fonction
$f\colon\mathbb{N}^k\dasharrow\mathbb{N}$ est générale récursive si et
seulement si elle est représentable par un terme du $\lambda$-calcul
non typé.
\end{thm}

\begin{proof}
Toute fonction générale récursive est représentable~: les fonctions
initiales et la composition sont directement représentables, la
récursion primitive l'est à l'aide du codage des couples, et la
minimisation l'est à l'aide d'un combinateur de point fixe.

Réciproquement, on code les termes par des entiers.  Les étapes de
réduction extérieure gauche sont primitives récursives et une
recherche non bornée trouve la première forme normale.  La lecture
d'un entier de Church dans cette forme normale est elle aussi
effective.  La fonction représentée est donc générale récursive.
\end{proof}

Cette équivalence est constructive~: on peut traduire
algorithmiquement un indice de fonction récursive en un
$\lambda$-terme qui la représente, et réciproquement.

\subsection{Conclusion}\label{subsection-computability-conclusion}

\thingy Nous avons obtenu trois descriptions équivalentes des fonctions
calculables.  Pour une fonction partielle
\[
f\colon\mathbb{N}^k\dasharrow\mathbb{N},
\]
les propriétés suivantes sont équivalentes~:
\begin{enumerate}
\item $f$ est générale récursive~;
\item $f$ est représentable dans le $\lambda$-calcul non typé~;
\item $f$ est calculable par une machine de Turing.
\end{enumerate}
Ces fonctions sont les \defin{fonctions calculables} au sens de
Church-Turing.

Les équivalences sont constructives~: elles fournissent des traductions
algorithmiques entre les différentes représentations.  On peut les
comprendre comme des compilateurs transformant un programme écrit dans
un modèle en un programme équivalent écrit dans un autre.

\thingy Ces modèles ont également les mêmes limitations.  Il n'existe
pas d'algorithme permettant de déterminer si une machine de Turing
donnée s'arrête, ni de déterminer si la réduction d'un
$\lambda$-terme donné atteint une forme normale.  Plus généralement,
le théorème de Rice montre qu'aucune propriété non triviale de la
fonction calculée par un programme n'est décidable à partir du code de
ce programme.

\thingy Un langage de programmation est dit
\defin[Turing-complet (langage)]{Turing-complet} lorsque, après une
idéalisation convenable, il permet de calculer exactement les fonctions
calculables au sens de Church-Turing.  Les langages généralistes usuels
sont Turing-complets~: Python, Java, JavaScript, C, C++, OCaml, Haskell,
Lisp, Perl, Ruby, Smalltalk ou Prolog, par exemple.

Certains langages ou systèmes le sont plus ou moins accidentellement,
par exemple certains fragments ou usages de CSS, \TeX, XSLT ou
\texttt{gm}.  La Turing-complétude ne signifie cependant pas qu'un
langage peut commodément exprimer n'importe quelle opération.  Elle
signifie seulement qu'il peut, en principe et avec les codages
appropriés, calculer toute fonction calculable.  La traduction peut
être extrêmement inefficace, illisible ou malcommode.

Un ordinateur réel ne peut pas calculer davantage qu'une machine de
Turing, sauf à disposer d'une véritable source de hasard qui modifierait
la nature des résultats attendus.  La question pertinente est plutôt
de savoir si un langage donné permet de calculer \emph{autant} qu'une
machine de Turing.

\thingy Enfin, le typage ne garantit pas à lui seul la terminaison.
L'existence d'un type $t$ tel que
\[
t\cong(t\to t)
\]
permet d'encoder dans le type $t$ les termes du $\lambda$-calcul non
typé.  On peut ainsi provoquer une boucle infinie dans un langage
apparemment typé, sans employer explicitement une définition
récursive.

\review{La construction OCaml donnée dans les diapositives utilise un
type récursif pour représenter l'auto-application.  Elle pourra être
reprise comme exemple ou exercice, mais elle n'est pas reproduite ici
afin de respecter la consigne de ne pas intégrer encore les
exercices.}

\thingy Les fonctions calculables peuvent croître extraordinairement
vite.  La fonction d'Ackermann en est un premier exemple, tandis que la
fonction du castor affairé domine toute fonction calculable.  Il est
donc possible de concevoir un programme de taille raisonnable qui
termine théoriquement, mais dont le temps d'exécution ou le résultat
dépasse toute grandeur utilisable en pratique.

La calculabilité distingue ainsi deux questions radicalement
différentes~:
\begin{itemize}
\item un calcul peut-il être effectué par un algorithme~?
\item s'il le peut, peut-il être effectué avec des ressources
  raisonnables~?
\end{itemize}
La première relève de la calculabilité~; la seconde, de la complexité.

%% PREAMBLE: Le chapitre utilise \defin, dont la définition appelle
%% \index.  Si ce n'est pas déjà fait dans le document maître, ajouter
%% \usepackage{makeidx} puis \makeindex dans le préambule.
%% PREAMBLE: Pour compiler l'éventuel index, utiliser par exemple
%% texindy -C utf8 -L french notes-inf110.idx.

\section{Généralités sur le typage}

\review{Ce chapitre constitue un premier brouillon produit avec l'aide
d'un modèle d'IA. Il doit être relu et vérifié par l'enseignant,
notamment en ce qui concerne les formulations mathématiques, les
preuves esquissées et les choix éditoriaux.}

\subsection{Rôle et variétés des systèmes de typage}

\thingy Informellement, un \defin{système de typage} est une façon
d'affecter aux expressions et aux valeurs manipulées par un langage
informatique un \defin{type} qui contraint leur usage.  Parmi les types
les plus courants, on trouve par exemple les entiers, les booléens, les
chaînes de caractères et les fonctions.

Cette manière de distinguer différentes sortes de données est, dans
une certaine mesure, opposée à la philosophie du codage de Gödel
rencontrée en calculabilité : dans ce dernier, toute donnée finie est
représentée par un entier, tandis qu'un système de typage cherche au
contraire à conserver et exploiter les distinctions entre les
différentes natures de données.

Les buts d'un système de typage sont multiples :
\begin{itemize}
\item détecter le plus tôt possible certaines erreurs de
  programmation : ajouter un entier et une chaîne de caractères est
  probablement une erreur, ou demande au moins une conversion
  explicite ;
\item éviter des plantages ou des problèmes de sécurité, par exemple
  lorsqu'un programme tente d'exécuter comme du code une donnée qui
  n'est qu'un entier ;
\item garantir certaines propriétés du comportement des programmes,
  dans les limites imposées par les résultats d'indécidabilité tels
  que le théorème de Rice ;
\item fournir des informations utiles à l'optimisation, par exemple en
  indiquant qu'une fonction est pure, c'est-à-dire sans effet de bord.
\end{itemize}
L'analogie avec les systèmes d'unités en physique est utile : de même
que l'homogénéité interdit d'ajouter directement une longueur à une
durée, le typage interdit certaines combinaisons d'expressions qui
n'ont pas de sens dans le langage considéré.

\thingy Il existe presque autant de systèmes de typage que de langages
de programmation. Plusieurs distinctions sont particulièrement
importantes.

Un typage est dit \defin{statique} lorsqu'il est vérifié au moment de la
compilation, et \defin{dynamique} lorsqu'il est vérifié pendant
l'exécution. Certains langages combinent les deux approches ; on parle
notamment de typage graduel. Du point de vue du programmeur, on préfère
en général une erreur de compilation à une erreur à l'exécution, et une
erreur à l'exécution à un plantage ou à un comportement indéfini.

Le typage peut être \defin{annoté}, lorsque le programmeur écrit
explicitement les types attendus, ou \defin{inféré}, lorsque le langage
les détermine à partir de l'expression. Cette distinction peut
s'interpréter de la manière suivante : le type est-il une promesse
faite par le langage au programmeur, ou une promesse faite par le
programmeur au langage ?

Il faut également distinguer les systèmes de typage dont les garanties
sont difficiles à contourner de ceux qui permettent des conversions
explicites (\textit{casts}), la manipulation directe des
représentations en mémoire, ou le report de certaines vérifications à
l'exécution.

Enfin, le typage peut être superficiel, par exemple en affirmant
seulement qu'une valeur est une liste, ou plus précis, en affirmant
qu'elle est une liste d'entiers. Il peut aussi enregistrer des
informations qui ne décrivent pas directement la représentation d'une
valeur : pureté d'une fonction, exceptions qu'elle peut soulever,
mutabilité, droits d'accès, utilisation de la mémoire, etc.

\thingy Dans certains langages, les types eux-mêmes sont des objets
manipulables par les programmes : on dit alors qu'ils sont des
\defin{citoyens de première classe}. Cette possibilité pose aussitôt la
question du type des types eux-mêmes. Elle conduit à des systèmes très
expressifs, mais aussi à des difficultés logiques et algorithmiques
importantes.

\subsection{Constructions usuelles sur les types}

\thingy En plus de types de base tels que $\texttt{Nat}$, pour les
entiers naturels, ou $\texttt{Bool}$, pour les booléens, les systèmes de
typage proposent souvent plusieurs opérations de construction de
types.

Les \defin{types produits} représentent les couples, triplets et, plus
généralement, les $k$-uplets. Par exemple,
\[
\texttt{Nat}\times\texttt{Bool}
\]
est le type des paires dont la première composante est un entier et la
seconde un booléen. Lorsque les composantes sont nommées, on obtient
des structures, aussi appelées enregistrements. Le produit vide est le
type trivial $\texttt{Unit}$, qui ne possède qu'une seule valeur.

Les \defin{types sommes} représentent un choix entre plusieurs cas. Le
type
\[
\texttt{Nat}+\texttt{Bool}
\]
contient soit un entier, soit un booléen, avec un sélecteur indiquant
lequel des deux cas est utilisé. Une somme de plusieurs copies du type
$\texttt{Unit}$ donne un type d'énumération, dont la seule information
est précisément ce sélecteur. La somme vide est un type inhabité,
c'est-à-dire un type ne possédant aucune valeur.

Les \defin{types fonctions}, également appelés types exponentiels, sont
notés à l'aide d'une flèche. Ainsi,
\[
\texttt{Nat}\rightarrow\texttt{Bool}
\]
est le type des fonctions de $\texttt{Nat}$ vers $\texttt{Bool}$.

Enfin, on rencontre fréquemment des types de listes. Par exemple,
$\texttt{List}~\texttt{Nat}$ est le type des listes d'entiers.

\thingy Plusieurs fonctionnalités peuvent enrichir ces constructions.

Le \defin{sous-typage} signifie que toute valeur d'un certain type peut
automatiquement être utilisée comme une valeur d'un autre type.

Le \defin{polymorphisme} permet à une expression d'être utilisée avec
plusieurs types. L'exemple fondamental est la fonction identité, dont
on peut écrire informellement le type
\[
(\forall\mathtt{t})\; \mathtt{t}\rightarrow\mathtt{t}.
\]

Une \defin{famille de types} fabrique un type à partir d'un ou plusieurs
autres types. Ainsi, $\texttt{List}$ transforme le type $\mathtt{t}$ en
le type $\texttt{List}~\mathtt{t}$ des listes d'éléments de type
$\mathtt{t}$.

Un \defin{type récursif} est construit, à l'aide des opérations sur les
types, en faisant intervenir le type défini lui-même. Les diapositives
donnent comme exemple schématique
\[
\texttt{Tree}=\texttt{List}~\texttt{Tree}.
\]
\review{Cet exemple décrit des arbres dont chaque nœud est essentiellement
une liste de sous-arbres, mais ne fournit pas de cas de base explicite ;
il convient de vérifier si cette présentation volontairement
schématique doit être conservée.}

Un \defin{type dépendant} est un type qui dépend d'une valeur. Par
exemple, à un entier $k$ on peut associer le type $\texttt{Nat}^k$ des
$k$-uplets d'entiers.

Enfin, un \defin{type opaque}, abstrait ou privé, est un type dont la
représentation interne est cachée. Ses valeurs ne peuvent alors être
manipulées qu'au moyen d'une interface publique.

\subsection{Polymorphisme et tâches d'un système de typage}

\thingy On distingue principalement deux sortes de polymorphisme.

Le \defin{polymorphisme paramétrique}, ou polymorphisme générique,
signifie que la même expression s'applique de la même manière à des
données de types différents. Quelques exemples sont :
\[
\texttt{head} :
(\forall\mathtt{t})\;
\texttt{List}~\mathtt{t}\to\mathtt{t},
\]
\[
\lambda xy.\langle x,y\rangle :
(\forall\mathtt{u},\mathtt{v})\;
\mathtt{u}\to\mathtt{v}\to
\mathtt{u}\times\mathtt{v},
\]
et
\[
\lambda xyz.xz(yz) :
(\forall\mathtt{u},\mathtt{v},\mathtt{w})\;
(\mathtt{u}\to\mathtt{v}\to\mathtt{w})
\to(\mathtt{u}\to\mathtt{v})
\to\mathtt{u}\to\mathtt{w}.
\]
Le polymorphisme ne concerne pas seulement les fonctions. La liste
vide, par exemple, a informellement le type
\[
\texttt{[]} :
(\forall\mathtt{t})\;\texttt{List}~\mathtt{t}.
\]
Dans un langage autorisant les boucles infinies, une expression qui ne
termine jamais peut également recevoir n'importe quel type, puisqu'elle
ne produit jamais de valeur contredisant ce type.

Le \defin{polymorphisme ad hoc}, ou surcharge, désigne au contraire une
opération qui agit différemment selon le type de son argument, ce type
étant connu au moment où l'opération est choisie.

Les frontières terminologiques sont parfois floues : le sous-typage,
et quelquefois même les coercions, sont présentés comme des formes de
polymorphisme.

\thingy Un système de typage peut accomplir plusieurs tâches.

La \defin{vérification de type} consiste à vérifier qu'une expression
annotée possède bien le type annoncé par ses annotations.

L'\defin{inférence de type}, également appelée reconstruction ou
assignation de type, consiste à déterminer le type d'une expression en
l'absence d'annotations, ou en présence d'annotations seulement
partielles. L'algorithme de Hindley--Milner, étudié à la fin de ce
chapitre, est un algorithme fondamental d'inférence pour les langages
fonctionnels à polymorphisme paramétrique.

Dans les systèmes suffisamment complexes, l'inférence, et parfois même
la vérification, peuvent devenir indécidables. Cela se produit
notamment lorsque les types sont des objets de première classe ou
peuvent dépendre de valeurs arbitraires calculées pendant l'exécution.

Enfin, le problème de l'\defin{habitation d'un type} consiste à trouver
un terme ayant un type donné. Cette question est essentielle en
logique, où elle correspondra à la recherche d'une preuve. Par exemple,
on peut demander s'il existe un terme de type
\[
(\forall\mathtt{p},\mathtt{q})\;
((\mathtt{p}\to\mathtt{q})\to\mathtt{p})\to\mathtt{p}.
\]
Dans les langages usuels qui autorisent la non-terminaison, tout type
peut être artificiellement habité par une boucle infinie, y compris un
type censé être vide. Cette observation explique pourquoi la
terminaison sera indispensable lorsque les types seront interprétés
comme des propositions.

\subsection{Autres usages et exemples}

\thingy Les types peuvent annoter des propriétés qui ne sont pas
simplement la nature des valeurs stockées. Ils peuvent indiquer les
exceptions susceptibles d'être soulevées, comme en Java, ou distinguer
les valeurs immuables des références mutables. Ainsi,
$\texttt{Nat}$ peut désigner un entier immuable, tandis que
$\texttt{Ref~Nat}$ désigne une référence mutable vers un entier.

Ils peuvent également enregistrer les effets de bord. En Haskell,
$\texttt{Char}$ représente une valeur pure de type caractère, tandis
que $\texttt{IO~Char}$ représente une action ayant des effets de bord
et renvoyant un caractère.

Le \defin{typage linéaire}, ou les types à unicité, imposent qu'une
valeur soit utilisée une et une seule fois au cours d'un calcul, sauf
opération spéciale. Ils interdisent donc aussi bien sa duplication que
sa perte. Cette information peut servir à la gestion de la mémoire,
comme en Rust, ou à l'annotation des effets de bord, comme en Clean.

\thingy Donnons quelques exemples de systèmes de typage, sans employer
les expressions ambiguës de typage « faible » ou « fort ».

L'assembleur ne possède essentiellement aucun typage : toute donnée
est représentée par une suite de bits. La même remarque vaut, dans des
cadres théoriques différents, pour les machines de Turing, les
fonctions générales récursives et le $\lambda$-calcul non typé.

Le langage C possède un typage annoté et vérifié à la compilation, mais
celui-ci est contournable, notamment à l'aide des pointeurs génériques
de type \texttt{void*}. Python, JavaScript, Perl ou Scheme effectuent
principalement les vérifications à l'exécution. Java combine des
vérifications à la compilation et à l'exécution ; les types génériques
ont été ajoutés avec Java 5. Rust fait interagir le typage et la gestion
de la mémoire au moyen de mécanismes proches du typage affine ou
linéaire.

Les langages fonctionnels ont souvent des systèmes plus complexes.
OCaml dispose notamment de l'inférence de type à la Hindley--Milner, du
polymorphisme paramétrique, de types récursifs et d'un système de
modules. Haskell possède des fonctionnalités comparables, auxquelles
s'ajoutent les classes de types et l'encapsulation des effets de bord
dans des monades. Mercury combine des idées issues de Haskell et de
Prolog. Idris est à la fois un langage fonctionnel et un assistant de
preuve.

Parmi les exemples plus spécialisés, F\# peut intégrer les unités de
mesure au système de typage, Rust peut exclure certaines situations de
concurrence dangereuses, et divers langages garantissent des propriétés
telles que des bornes de complexité, la validation de documents XML ou
l'absence de fuite d'information.

\subsection{Typage, terminaison et correspondance de Curry--Howard}

\thingy\label{typing-and-termination} Un système de typage peut garantir
que tout programme bien typé termine. Ce n'est pas le cas des langages
usuels tels qu'OCaml ou Haskell, dans lesquels un programme bien typé
peut effectuer une boucle infinie.

Si le typage doit lui-même être décidable, garantir la terminaison
impose nécessairement des limites à l'expressivité du langage, en
raison de l'indécidabilité du problème de l'arrêt. Un tel langage ne
peut notamment pas écrire son propre interpréteur universel. Il ne
peut pas non plus autoriser des boucles illimitées, des appels récursifs
non contraints ou un type $\mathtt{t}$ satisfaisant
\[
\mathtt{t}\cong(\mathtt{t}\rightarrow\mathtt{t}).
\]
Son type vide doit réellement être inhabité.

Ces restrictions n'empêchent cependant pas le langage de définir des
fonctions beaucoup plus générales que les seules fonctions primitives
récursives. Rocq, anciennement appelé Coq, fournit un exemple de
langage et d'assistant de preuve garantissant la terminaison.

\thingy La \defin{correspondance de Curry--Howard} établit une analogie,
et dans des systèmes précisément définis une correspondance exacte,
entre :
\begin{itemize}
\item les types et leurs termes, c'est-à-dire les programmes possédant
  ces types ;
\item les propositions et leurs preuves.
\end{itemize}
Par exemple, la preuve évidente de
\[
(A\land B)\Rightarrow A
\]
correspond à la première projection, de type
\[
a\times b\rightarrow a.
\]
De même, le terme $\lambda xy.x$, de type $u\to v\to u$, correspond à
la preuve de la proposition
\[
U\Rightarrow V\Rightarrow U.
\]

Cette correspondance exige certaines précautions. Le langage de
programmation doit garantir la terminaison, faute de quoi une boucle
infinie jouerait le rôle d'une preuve circulaire. La logique
correspondante est en général intuitionniste, donc dépourvue du tiers
exclu. Enfin, les quantificateurs et les types dépendants font
intervenir des subtilités supplémentaires.

Dans un premier aperçu, la correspondance associe :
\begin{itemize}
\item le type fonction $A\to B$ à l'implication $A\Rightarrow B$ ;
\item l'application d'une fonction au \textit{modus ponens} ;
\item l'abstraction à l'ouverture puis à la décharge d'une hypothèse ;
\item le produit $A\times B$ à la conjonction $A\land B$ ;
\item la somme $A+B$ à la disjonction $A\lor B$ ;
\item les types trivial et vide à $\top$ et $\bot$.
\end{itemize}
La négation demandera une discussion particulière.

\thingy Le paradoxe de Curry illustre le danger des preuves
circulaires. Supposons que l'on puisse construire une proposition $A$
telle que
\[
A\Leftrightarrow(A\Rightarrow B).
\]
Alors $B$ est effectivement démontrable : sous l'hypothèse $A$,
l'équivalence donne $A\Rightarrow B$, donc $B$. Cela établit
$A\Rightarrow B$, qui, par l'équivalence, donne à son tour $A$, puis
finalement $B$.

L'erreur ne se trouve donc pas dans cette déduction, mais dans la
supposition tacite qu'une proposition $A$ satisfaisant une telle
auto-référence peut toujours être construite. Une construction de
Quine permet une auto-référence syntaxique, mais ne suffit pas à
obtenir une auto-référence sur la vérité.

\thingy Les \defin{théories des types} utilisées en logique sont des
langages permettant d'écrire à la fois des programmes et des preuves.
Elles peuvent être étudiées comme des systèmes abstraits ou
implémentées dans des assistants de preuve tels que Rocq, Agda ou Lean.
Elles garantissent la terminaison des programmes, ou plus précisément,
dans le cadre du $\lambda$-calcul, la normalisation forte.

Parmi les grandes familles, on peut citer les théories des types à la
Martin-Löf, les variantes du calcul des constructions, et la théorie
homotopique des types. Le $\lambda$-cube de Barendregt organise huit
théories selon la présence ou l'absence de trois fonctionnalités :
\begin{itemize}
\item les termes peuvent dépendre de types ;
\item les types peuvent dépendre de types ;
\item les types peuvent dépendre de termes.
\end{itemize}
Le $\lambda$-calcul simplement typé ne possède aucune de ces trois
fonctionnalités, tandis que le calcul des constructions les possède
toutes.

\section{Le \texorpdfstring{$\lambda$}{lambda}-calcul simplement typé}

\subsection{Syntaxe et règles de typage}

\thingy Le \defin{$\lambda$-calcul simplement typé}, abrégé
$\lambda$CST et également noté $\lambda_\rightarrow$, est une variante
typée du $\lambda$-calcul. Il ne possède qu'une seule opération sur les
types : si $\sigma$ et $\tau$ sont deux types, alors
$\sigma\to\tau$ est le type des fonctions qui prennent un argument de
type $\sigma$ et renvoient une valeur de type $\tau$.

L'application et l'abstraction doivent respecter les types :
\begin{itemize}
\item si $P$ a pour type $\sigma\to\tau$ et $Q$ a pour type $\sigma$,
  alors $PQ$ a pour type $\tau$ ;
\item si $E$ a pour type $\tau$ sous l'hypothèse que la variable libre
  $v$ a pour type $\sigma$, alors
  $\lambda(v:\sigma).E$ a pour type $\sigma\to\tau$.
\end{itemize}

Le typage considéré est annoté, ou « à la Church » : on écrit
$\lambda(v:\sigma).E$, et non simplement $\lambda v.E$.

\thingy Un \defin{type} du $\lambda$CST est défini inductivement par les
deux clauses suivantes :
\begin{itemize}
\item toute variable de type $\alpha,\beta,\gamma,\ldots$ est un type ;
\item si $\sigma$ et $\tau$ sont des types, alors
  $(\sigma\to\tau)$ est un type.
\end{itemize}

Un \defin{préterme} est défini inductivement par :
\begin{itemize}
\item toute variable de terme est un préterme ;
\item si $P$ et $Q$ sont des prétermes, alors $(PQ)$ est un préterme ;
\item si $v$ est une variable, $\sigma$ un type et $E$ un préterme,
  alors $\lambda(v:\sigma).E$ est un préterme.
\end{itemize}

On adopte les conventions de parenthésage
\[
\rho\to\sigma\to\tau
  :=\rho\to(\sigma\to\tau)
\qquad\text{et}\qquad
xyz:=((xy)z).
\]
L'application est donc associative à gauche dans la notation, tandis
que le constructeur de type fonction est associatif à droite.
L'abstraction est moins prioritaire que l'application. On peut abréger
\[
\lambda(x:\sigma).\lambda(t:\tau).E
\]
en $\lambda(x:\sigma,t:\tau).E$ ou
$\lambda(x:\sigma)(t:\tau).E$. Enfin, les termes sont considérés à
renommage près de leurs variables liées, c'est-à-dire à
$\alpha$-conversion près.

\begin{defn}
Un \defin{contexte de typage} est un ensemble fini $\Gamma$ de typages
$x_i:\sigma_i$, où les $x_i$ sont des variables de terme distinctes.
On l'écrit généralement
\[
\Gamma=x_1:\sigma_1,\ldots,x_k:\sigma_k.
\]

Un \defin{jugement de typage} est une expression
\[
\Gamma\vdash E:\tau,
\]
signifiant que $E$ est un terme de type $\tau$ dans le contexte
$\Gamma$. Lorsque $\Gamma$ est vide, on écrit simplement
$\vdash E:\tau$.
\end{defn}

\thingy Les règles de typage du $\lambda$CST sont les suivantes :
\[
\inferrule*[Left={Var}]
 { }
 {\Gamma,x:\sigma\vdash x:\sigma}
\]
\[
\inferrule*[Left={App},sep=4em]
 {\Gamma\vdash P:\sigma\to\tau\\
  \Gamma\vdash Q:\sigma}
 {\Gamma\vdash PQ:\tau}
\]
\[
\inferrule*[Left={Abs}]
 {\Gamma,v:\sigma\vdash E:\tau}
 {\Gamma\vdash\lambda(v:\sigma).E:\sigma\to\tau}.
\]
Chaque barre horizontale signifie que le jugement placé en dessous
découle, par la règle indiquée, des jugements placés au-dessus.
L'ensemble des jugements de typage est le plus petit ensemble fermé
par ces règles.

Une \defin{dérivation} d'un jugement est un arbre d'instances de ces
règles dont la racine est le jugement considéré.

\thingy Voici quelques exemples de jugements :
\[
f:\alpha\to\beta,\;x:\alpha\vdash fx:\beta,
\]
\[
f:\beta\to\alpha\to\gamma,\;x:\alpha,\;y:\beta
  \vdash fyx:\gamma,
\]
\[
f:\beta\to\alpha\to\gamma,\;x:\alpha
  \vdash\lambda(y:\beta).fyx:\beta\to\gamma,
\]
\[
x:\alpha\vdash
  \lambda(f:\alpha\to\beta).fx:(\alpha\to\beta)\to\beta,
\]
et
\[
\vdash
\lambda(x:\alpha).\lambda(f:\alpha\to\beta).fx:
\alpha\to(\alpha\to\beta)\to\beta.
\]

À titre d'exemple plus développé, on obtient la dérivation suivante :
{\footnotesize
\[
\inferrule*[Left={Abs}]{
 \inferrule*[Left={Abs}]{
  \inferrule*[Left={Abs}]{
   \inferrule*[Left={App}]{
    \inferrule*[Left={App}]{
     \inferrule*[Left={Var}]{ }
      {f:\beta\to\alpha\to\gamma,x:\alpha,y:\beta
       \vdash f:\beta\to\alpha\to\gamma}
     \\
     \inferrule*[Right={Var}]{ }
      {f:\beta\to\alpha\to\gamma,x:\alpha,y:\beta
       \vdash y:\beta}
    }{
     f:\beta\to\alpha\to\gamma,x:\alpha,y:\beta
     \vdash fy:\alpha\to\gamma
    }
    \\
    \inferrule*[Right={Var}]{ }
     {f:\beta\to\alpha\to\gamma,x:\alpha,y:\beta
      \vdash x:\alpha}
   }{
    f:\beta\to\alpha\to\gamma,x:\alpha,y:\beta
    \vdash fyx:\gamma
   }
  }{
   f:\beta\to\alpha\to\gamma,x:\alpha
   \vdash\lambda(y:\beta).fyx:\beta\to\gamma
  }
 }{
  f:\beta\to\alpha\to\gamma
  \vdash\lambda(x:\alpha).\lambda(y:\beta).fyx:
  \alpha\to\beta\to\gamma
 }
}{
 \vdash
 \lambda(f:\beta\to\alpha\to\gamma).
 \lambda(x:\alpha).\lambda(y:\beta).fyx:
 (\beta\to\alpha\to\gamma)\to(\alpha\to\beta\to\gamma)
}
\]
}

\subsection{Propriétés élémentaires du typage}

\begin{prop}[Affaiblissement]
Si $\Gamma\subseteq\Gamma'$ sont deux contextes et si
$\Gamma\vdash M:\sigma$, alors $\Gamma'\vdash M:\sigma$.
\end{prop}

\begin{proof}
On raisonne par induction sur la dérivation de
$\Gamma\vdash M:\sigma$. Dans le cas d'une variable, le typage utilisé
appartient à $\Gamma$, donc également à $\Gamma'$. Les règles
d'application et d'abstraction se transportent immédiatement en
appliquant l'hypothèse d'induction à leurs prémisses.
\end{proof}

\begin{prop}[Variables libres]
Si
\[
x_1:\sigma_1,\ldots,x_k:\sigma_k\vdash E:\tau,
\]
alors toute variable libre de $E$ appartient à
$\{x_1,\ldots,x_k\}$.
\end{prop}

\begin{proof}
La propriété se vérifie par induction sur la dérivation. La règle
variable ne peut introduire qu'une variable présente dans le contexte.
L'application réunit les variables libres de ses deux sous-termes.
Enfin, l'abstraction lie précisément la variable ajoutée au contexte.
\end{proof}

En conséquence, les typages du contexte qui ne concernent aucune
variable libre de $E$ peuvent être supprimés. Le typage est également
invariant par renommage cohérent des variables libres.

\thingy La dérivation du typage d'un préterme se construit
systématiquement à partir de sa syntaxe :
\begin{itemize}
\item une variable doit être présente dans le contexte ;
\item pour une application $PQ$, les types obtenus pour $P$ et $Q$
  doivent être respectivement $\sigma\to\tau$ et $\sigma$ ;
\item pour une abstraction $\lambda(v:\sigma).E$, on type $E$ dans le
  contexte enrichi par $v:\sigma$.
\end{itemize}
Il n'y a donc aucun choix substantiel à effectuer.

\begin{prop}[Unicité du type]
Pour un contexte $\Gamma$ et un préterme annoté $M$ donnés, il existe
au plus un type $\sigma$ tel que $\Gamma\vdash M:\sigma$.
\end{prop}

\begin{proof}
On raisonne par induction sur la structure de $M$. Le type d'une
variable est imposé par le contexte. Celui d'une abstraction est imposé
par l'annotation de sa variable et par le type, unique par hypothèse
d'induction, de son corps. Enfin, dans une application, le type de la
fonction et celui de l'argument imposent de manière unique le type du
résultat.
\end{proof}

\subsection{Substitution, réduction et normalisation}

\begin{lem}[Substitution]\label{substitution-lemma}
Si
\[
\Gamma,v:\sigma\vdash E:\tau
\qquad\text{et}\qquad
\Gamma\vdash T:\sigma,
\]
alors
\[
\Gamma\vdash E[v\backslash T]:\tau,
\]
où $E[v\backslash T]$ désigne la substitution correcte de $v$ par $T$
dans $E$.
\end{lem}

\begin{proof}
On raisonne par induction sur la dérivation de
$\Gamma,v:\sigma\vdash E:\tau$. Les règles d'application et
d'abstraction se transportent en appliquant l'hypothèse d'induction aux
prémisses, après avoir renommé si nécessaire les variables liées pour
éviter toute capture.

Dans le cas de la règle variable, si la variable considérée est $v$,
on remplace la feuille correspondante par la dérivation de
$\Gamma\vdash T:\sigma$. Si elle est différente de $v$, la même règle
variable reste applicable dans $\Gamma$. Ceci construit une dérivation
de $\Gamma\vdash E[v\backslash T]:\tau$.
\end{proof}

\thingy La $\beta$-réduction remplace un redex
\[
(\lambda(v:\sigma).E)T
\]
par son réduit
\[
E[v\backslash T].
\]
On note $X\rightarrow_\beta X'$ pour une étape de
$\beta$-réduction et $X\twoheadrightarrow_\beta X'$ pour une suite de
telles étapes.

\begin{prop}[Préservation du typage]
Si $\Gamma\vdash X:\tau$ et
$X\twoheadrightarrow_\beta X'$, alors
$\Gamma\vdash X':\tau$.
\end{prop}

\begin{proof}
Il suffit de considérer une étape de réduction. Lorsqu'un redex
$(\lambda(v:\sigma).E)T$ est bien typé, on dispose de
\[
\Gamma,v:\sigma\vdash E:\tau
\qquad\text{et}\qquad
\Gamma\vdash T:\sigma.
\]
Le lemme de substitution donne alors
$\Gamma\vdash E[v\backslash T]:\tau$. La réduction à l'intérieur d'un
contexte se traite par induction sur ce contexte. Une suite de
réductions se traite ensuite par induction sur sa longueur.
\end{proof}

\begin{thm}[Turing, Tait, Girard]
Le $\lambda$-calcul simplement typé est \defin{fortement normalisant} :
toute suite de $\beta$-réductions issue d'un terme bien typé est finie.
\end{thm}

\begin{proof}
Les diapositives indiquent la méthode dite de calculabilité. On définit
par induction sur le type la notion de terme \emph{fortement
calculable} :
\begin{itemize}
\item un terme de type atomique est fortement calculable lorsqu'il est
  fortement normalisable ;
\item un terme $P$ de type $\sigma\to\tau$ est fortement calculable
  lorsque, pour tout terme $Q$ de type $\sigma$ fortement calculable,
  le terme $PQ$ est fortement calculable.
\end{itemize}
On montre d'abord, par induction sur le type, que tout terme fortement
calculable est fortement normalisable. On montre ensuite, par induction
sur une dérivation de typage, que la substitution de termes fortement
calculables aux variables libres d'un terme bien typé produit un terme
fortement calculable. En substituant les variables elles-mêmes, on
obtient que tout terme bien typé est fortement calculable, donc
fortement normalisable.

\review{La preuve précédente développe seulement l'esquisse donnée dans
les diapositives. Une version finale devrait préciser les propriétés de
stabilité des termes fortement calculables utilisées dans le cas de
l'abstraction.}
\end{proof}

\subsection{Reconstruction des termes et problèmes algorithmiques}

\thingy La syntaxe du terme détermine directement la forme de sa
dérivation de typage. Réciproquement, si l'on connaît un arbre de
dérivation dans lequel les types et les règles sont indiqués, mais où
les termes ont été effacés, les termes peuvent être reconstruits en
suivant l'arbre :
\begin{itemize}
\item une règle variable détermine la variable utilisée ;
\item une règle d'application reconstruit l'application des deux termes
  obtenus dans les prémisses ;
\item une règle d'abstraction reconstruit la variable liée et le corps.
\end{itemize}
Il faut seulement distinguer correctement les variables lorsque
plusieurs hypothèses du contexte ont le même type.

\thingy Les problèmes suivants ont des difficultés différentes :
\begin{itemize}
\item étant donné un préterme annoté, calculer son type ou vérifier un
  jugement est facile et décidable ;
\item étant donné un arbre de dérivation typé dont les termes ont été
  effacés, reconstruire le terme est également facile ;
\item étant donné un terme non annoté, déterminer s'il admet des
  annotations de type demande un véritable algorithme d'inférence ;
\item étant donné un type, déterminer s'il est habité correspond, par
  Curry--Howard, à la décidabilité du calcul propositionnel
  intuitionniste.
\end{itemize}

\section{Le calcul propositionnel intuitionniste}

\subsection{Formules et déduction naturelle}

\thingy Le \defin{calcul propositionnel} est une logique dont les
objets de base sont des propositions, sans variables d'individus ni
quantificateurs $\forall$ ou $\exists$. Ses connecteurs sont
\[
\Rightarrow,\qquad \land,\qquad \lor,\qquad \top,\qquad \bot.
\]
La négation est une abréviation :
\[
\neg P:=P\Rightarrow\bot.
\]

Une \defin{formule propositionnelle} est définie inductivement comme
une variable propositionnelle, ou comme l'application d'un connecteur
aux formules appropriées. On adopte les priorités usuelles : $\neg$
est plus prioritaire que $\land$, qui est plus prioritaire que $\lor$,
qui est plus prioritaire que $\Rightarrow$. L'implication associe à
droite.

Un \defin{séquent}
\[
P_1,\ldots,P_r\vdash Q
\]
signifie que $Q$ est démontrable sous les hypothèses
$P_1,\ldots,P_r$.

\thingy La logique classique admet le tiers exclu, sous l'une de ses
formes équivalentes. La logique intuitionniste étudiée ici ne l'admet
pas. Il ne faut pas l'interpréter comme une logique possédant une
troisième valeur de vérité : elle impose plutôt une notion plus
constructive de preuve.

Chaque connecteur possède des règles d'\defin{introduction}, permettant
de le démontrer, et des règles d'\defin{élimination}, permettant de
l'utiliser.

Les règles de la déduction naturelle intuitionniste sont :
\[
\inferrule*[Left={Ax}]
 { }
 {\Gamma,P\vdash P}
\]
\[
\inferrule*[Left={$\Rightarrow$Int}]
 {\Gamma,P\vdash Q}
 {\Gamma\vdash P\Rightarrow Q}
\qquad
\inferrule*[Left={$\Rightarrow$Élim}]
 {\Gamma\vdash P\Rightarrow Q\\
  \Gamma\vdash P}
 {\Gamma\vdash Q}
\]
\[
\inferrule*[Left={$\land$Int}]
 {\Gamma\vdash Q_1\\
  \Gamma\vdash Q_2}
 {\Gamma\vdash Q_1\land Q_2}
\]
\[
\inferrule*[Left={$\land$Élim$_1$}]
 {\Gamma\vdash Q_1\land Q_2}
 {\Gamma\vdash Q_1}
\qquad
\inferrule*[Left={$\land$Élim$_2$}]
 {\Gamma\vdash Q_1\land Q_2}
 {\Gamma\vdash Q_2}
\]
\[
\inferrule*[Left={$\lor$Int$_1$}]
 {\Gamma\vdash Q_1}
 {\Gamma\vdash Q_1\lor Q_2}
\qquad
\inferrule*[Left={$\lor$Int$_2$}]
 {\Gamma\vdash Q_2}
 {\Gamma\vdash Q_1\lor Q_2}
\]
\[
\inferrule*[Left={$\lor$Élim}]
 {\Gamma\vdash P_1\lor P_2\\
  \Gamma,P_1\vdash Q\\
  \Gamma,P_2\vdash Q}
 {\Gamma\vdash Q}
\]
\[
\inferrule*[Left={$\top$Int}]
 { }
 {\Gamma\vdash\top}
\qquad
\inferrule*[Left={$\bot$Élim}]
 {\Gamma\vdash\bot}
 {\Gamma\vdash Q}.
\]

Dans ces règles, les hypothèses mises en évidence dans une prémisse
d'une règle d'introduction de $\Rightarrow$ ou d'élimination de $\lor$
sont déchargées dans la conclusion.

\thingy Par exemple, la commutativité de la conjonction se démontre par
\[
\inferrule*[Left={$\Rightarrow$Int}]{
 \inferrule*[Left={$\land$Int}]{
  \inferrule*[Left={$\land$Élim$_2$}]{
   \inferrule*[Left={Ax}]{ }{A\land B\vdash A\land B}
  }{A\land B\vdash B}
  \\
  \inferrule*[Right={$\land$Élim$_1$}]{
   \inferrule*[Right={Ax}]{ }{A\land B\vdash A\land B}
  }{A\land B\vdash A}
 }{A\land B\vdash B\land A}
}{\vdash A\land B\Rightarrow B\land A}.
\]

La commutativité de la disjonction se démontre en distinguant les deux
cas fournis par une preuve de $A\lor B$ :
\[
\inferrule*[Left={$\Rightarrow$Int}]{
 \inferrule*[Left={$\lor$Élim}]{
  \inferrule*[Left={Ax}]{ }{A\lor B\vdash A\lor B}
  \\
  \inferrule*[Right={$\lor$Int$_2$}]{
   \inferrule*[Right={Ax}]{ }{A\lor B,A\vdash A}
  }{A\lor B,A\vdash B\lor A}
  \\
  \inferrule*[Right={$\lor$Int$_1$}]{
   \inferrule*[Right={Ax}]{ }{A\lor B,B\vdash B}
  }{A\lor B,B\vdash B\lor A}
 }{A\lor B\vdash B\lor A}
}{\vdash A\lor B\Rightarrow B\lor A}.
\]

\subsection{Théorèmes et non-théorèmes}

\thingy Une formule $P$ telle que $\vdash P$ est dite
\defin{valable}, ou encore un \defin{théorème} du calcul
propositionnel intuitionniste.

Parmi les théorèmes importants figurent
\[
A\Rightarrow A,
\qquad
A\Rightarrow B\Rightarrow A,
\]
\[
(A\Rightarrow B\Rightarrow C)
\Rightarrow(A\Rightarrow B)\Rightarrow A\Rightarrow C,
\]
ainsi que les lois usuelles d'associativité, de commutativité et de
distributivité de $\land$ et $\lor$, sous leurs formes
intuitionnistement valables.

En revanche, les implications marquées dans les diapositives comme
classiques ne sont pas nécessairement démontrables
intuitionnistement. Parmi les non-théorèmes figurent :
\[
((A\Rightarrow B)\Rightarrow A)\Rightarrow A,
\]
\[
A\lor\neg A,
\qquad
\neg\neg A\Rightarrow A,
\]
\[
\neg(A\land B)\Rightarrow\neg A\lor\neg B,
\]
\[
(A\Rightarrow B)\Rightarrow\neg A\lor B,
\]
et
\[
(\neg B\Rightarrow\neg A)\Rightarrow(A\Rightarrow B).
\]
Ces formules sont valables en logique classique, mais pas dans la
logique intuitionniste.

\thingy Il faut distinguer une preuve de la négation d'une preuve par
l'absurde.

Pour démontrer $\neg A$, on suppose $A$ et on en déduit $\bot$. Cette
démonstration est intuitionnistement valable, puisqu'elle construit une
preuve de $A\Rightarrow\bot$.

Une preuve classique par l'absurde suppose au contraire $\neg A$,
obtient une contradiction, puis conclut $A$. Intuitionnistement, cet
argument démontre seulement $\neg\neg A$, et non nécessairement $A$.
Ainsi, $\neg A$ peut se lire « $A$ conduit à l'absurde », tandis que
$\neg\neg A$ signifie seulement que la négation de $A$ conduit à
l'absurde.

\subsection{Logique classique}

\thingy La logique classique peut être obtenue en remplaçant la règle
d'élimination de $\bot$ par une règle de raisonnement par l'absurde :
\[
\inferrule*[Left={Absurde}]
 {\Gamma,\neg Q\vdash\bot}
 {\Gamma\vdash Q}.
\]
Elle possède donc davantage de théorèmes que la logique
intuitionniste.

La sémantique booléenne fournit une caractérisation de la logique
classique : une formule $P$ est un théorème classique si et seulement
si toute affectation de $\bot$ ou $\top$ à ses variables lui donne la
valeur $\top$.

\begin{thm}[Correction et complétude classiques]
Les tables de vérité booléennes sont correctes et complètes pour le
calcul propositionnel classique.
\end{thm}

La correction se démontre par induction sur les dérivations. Pour la
complétude, on peut raisonner successivement sur les valeurs possibles
des variables et utiliser le tiers exclu afin de couvrir toutes les
lignes de la table de vérité.

\subsection{Interprétation de Brouwer--Heyting--Kolmogorov}

\thingy L'interprétation de Brouwer--Heyting--Kolmogorov donne une
lecture constructive des connecteurs :
\begin{itemize}
\item un témoignage de $P\land Q$ est un témoignage de $P$ accompagné
  d'un témoignage de $Q$ ;
\item un témoignage de $P\lor Q$ est un témoignage de l'un des deux,
  accompagné de l'indication du cas choisi ;
\item un témoignage de $P\Rightarrow Q$ est un moyen de transformer
  tout témoignage de $P$ en un témoignage de $Q$ ;
\item un témoignage de $\top$ est trivial ;
\item il n'existe aucun témoignage de $\bot$.
\end{itemize}
Cette interprétation explique directement les règles de déduction
naturelle.

\subsection{Correspondance de Curry--Howard}

\thingy Pour le connecteur $\Rightarrow$, les règles du $\lambda$CST et
celles de la déduction naturelle sont exactement les mêmes, au
changement de vocabulaire près :
\[
\begin{array}{c|c}
\text{Typage du $\lambda$CST}
&
\text{Calcul propositionnel intuitionniste}
\\[1ex]\hline\\[-2ex]
\inferrule*[Left={Var}]
 { }
 {\Gamma,x:\sigma\vdash x:\sigma}
&
\inferrule*[Left={Ax}]
 { }
 {\Gamma,P\vdash P}
\\[3ex]
\inferrule*[Left={App}]
 {\Gamma\vdash f:\sigma\to\tau\\
  \Gamma\vdash x:\sigma}
 {\Gamma\vdash fx:\tau}
&
\inferrule*[Left={$\Rightarrow$Élim}]
 {\Gamma\vdash P\Rightarrow Q\\
  \Gamma\vdash P}
 {\Gamma\vdash Q}
\\[3ex]
\inferrule*[Left={Abs}]
 {\Gamma,v:\sigma\vdash t:\tau}
 {\Gamma\vdash\lambda(v:\sigma).t:\sigma\to\tau}
&
\inferrule*[Left={$\Rightarrow$Int}]
 {\Gamma,P\vdash Q}
 {\Gamma\vdash P\Rightarrow Q}
\end{array}
\]
On peut donc identifier les termes du $\lambda$CST aux preuves en
déduction naturelle restreinte à l'implication.

Par exemple, le terme
\[
\lambda(f:\beta\to\alpha\to\gamma).
\lambda(x:\alpha).
\lambda(y:\beta).fyx
\]
correspond à une preuve de
\[
(B\Rightarrow A\Rightarrow C)
\Rightarrow(A\Rightarrow B\Rightarrow C).
\]

\thingy Pour traiter la conjonction, on enrichit le $\lambda$CST de
types produits. Les règles sont
\[
\inferrule*[Left={Pair}]
 {\Gamma\vdash t_1:\tau_1\\
  \Gamma\vdash t_2:\tau_2}
 {\Gamma\vdash\langle t_1,t_2\rangle:\tau_1\times\tau_2}
\]
et
\[
\inferrule*[Left={Proj$_i$}]
 {\Gamma\vdash t:\tau_1\times\tau_2}
 {\Gamma\vdash\pi_i t:\tau_i}
 \qquad(i\in\{1,2\}).
\]
Elles correspondent exactement aux règles d'introduction et
d'élimination de $\land$. Les réductions associées sont
\[
\pi_i\langle t_1,t_2\rangle\rightsquigarrow t_i.
\]
La preuve de $A\land B\Rightarrow B\land A$ correspond ainsi au terme
\[
\lambda(u:\alpha\times\beta).
\langle\pi_2u,\pi_1u\rangle.
\]

Pour traiter la disjonction, on ajoute des types sommes, avec deux
injections
\[
\iota_i^{(\tau_1,\tau_2)}:\tau_i\to\tau_1+\tau_2
\qquad(i\in\{1,2\}),
\]
et une construction par analyse de cas :
\[
\inferrule*[Left={Match}]
 {\Gamma\vdash r:\sigma_1+\sigma_2\\
  \Gamma,v_1:\sigma_1\vdash t_1:\tau\\
  \Gamma,v_2:\sigma_2\vdash t_2:\tau}
 {\Gamma\vdash
  (\texttt{match~}r\texttt{~with~}
   \iota_1v_1\mapsto t_1,\;
   \iota_2v_2\mapsto t_2):\tau}.
\]
Les réductions associées sont
\[
(\texttt{match~}\iota_i s\texttt{~with~}
 \iota_1v_1\mapsto t_1,\;
 \iota_2v_2\mapsto t_2)
\rightsquigarrow t_i[v_i\backslash s].
\]

Enfin, $\top$ correspond à un type unité $1$ possédant une unique
valeur $\bullet$, tandis que $\bot$ correspond à un type vide $0$.
Une élimination du type vide fournit un terme de n'importe quel type :
\[
\inferrule*[Left={Void}]
 {\Gamma\vdash r:0}
 {\Gamma\vdash\texttt{exfalso}^{(\tau)}r:\tau}.
\]

Ces extensions restent fortement normalisantes.

\thingy La correspondance de Curry--Howard peut alors se résumer ainsi :
\begin{center}
\begin{tabular}{c|c}
Types & Propositions\\ \hline
$\sigma\to\tau$ & $P\Rightarrow Q$\\
$\sigma\times\tau$ & $P\land Q$\\
$\sigma+\tau$ & $P\lor Q$\\
$1$ & $\top$\\
$0$ & $\bot$
\end{tabular}
\end{center}
Les termes correspondent aux preuves, l'abstraction à l'ouverture
d'une hypothèse, l'application au \textit{modus ponens}, et les
variables liées aux hypothèses déchargées.

Une même proposition peut avoir plusieurs preuves fonctionnellement
différentes. Par exemple, les entiers de Church
\[
\overline{0}_\alpha
 :=\lambda(f:\alpha\to\alpha).\lambda(x:\alpha).x,
\]
\[
\overline{1}_\alpha
 :=\lambda(f:\alpha\to\alpha).\lambda(x:\alpha).fx,
\]
\[
\overline{2}_\alpha
 :=\lambda(f:\alpha\to\alpha).\lambda(x:\alpha).f(fx)
\]
ont tous le type
\[
(\alpha\to\alpha)\to(\alpha\to\alpha)
\]
et correspondent donc à différentes preuves du même énoncé
$(A\Rightarrow A)\Rightarrow(A\Rightarrow A)$.

\subsection{Calcul des séquents et élimination des coupures}

\thingy Le \defin{calcul des séquents} est une autre présentation de la
même logique propositionnelle intuitionniste. Il distingue, pour
chaque connecteur, une règle droite qui l'introduit dans la conclusion
et une règle gauche qui l'introduit dans les hypothèses. On conserve
notamment l'axiome
\[
\inferrule*[Left={Ax}]{ }{\Gamma,Q\vdash Q},
\]
ainsi que les règles structurelles
\[
\inferrule*[Left={Contr}]
 {\Gamma,P,P\vdash Q}
 {\Gamma,P\vdash Q}
\qquad
\inferrule*[Left={Weak}]
 {\Gamma\vdash Q}
 {\Gamma,P\vdash Q}.
\]

La règle de \defin{coupure} est
\[
\inferrule*[Left={Cut}]
 {\Gamma\vdash P\\
  \Gamma,P\vdash Q}
 {\Gamma\vdash Q}.
\]
Elle exprime la composition de deux démonstrations, $P$ étant la
formule coupée.

\begin{thm}[Gentzen]
La règle de coupure est éliminable : tout séquent démontrable avec la
règle de coupure est démontrable sans elle.
\end{thm}

\begin{proof}
La preuve transforme localement une démonstration contenant une
coupure. On effectue une double récurrence : d'abord sur le degré de la
formule coupée, puis, à degré fixé, sur la somme des hauteurs des deux
branches de la coupure.

Lorsque la dernière règle de l'une des branches ne concerne pas la
formule coupée, on fait remonter la coupure au-dessus de cette règle.
Lorsqu'une branche se termine par l'axiome portant sur la formule
coupée, la coupure disparaît. Enfin, lorsque les deux dernières règles
introduisent le connecteur principal de la formule coupée, on remplace
la coupure par une ou plusieurs coupures portant sur des sous-formules
strictes. Le degré diminue alors, ce qui permet d'appliquer l'hypothèse
de récurrence.

\review{La preuve complète demande l'examen systématique de toutes les
paires de règles principales. Les diapositives n'en donnent que la
méthode et quelques exemples de transformations locales.}
\end{proof}

\thingy L'élimination des coupures a plusieurs conséquences.

Elle implique la \defin{propriété de la disjonction} :
\[
\text{si }\vdash Q_1\lor Q_2,\quad
\text{alors }\vdash Q_1\text{ ou }\vdash Q_2.
\]
En effet, une démonstration sans coupure d'une disjonction sans
hypothèse doit se terminer par l'une des deux règles d'introduction de
$\lor$.

Elle implique également la \defin{propriété de la sous-formule} : une
preuve sans coupure de
\[
P_1,\ldots,P_r\vdash Q
\]
ne fait intervenir que des sous-formules des $P_i$ et de $Q$.

On en déduit un algorithme de décision du calcul propositionnel
intuitionniste : il suffit d'énumérer les séquents construits avec les
sous-formules pertinentes et de fermer cet ensemble par les règles
autres que la coupure.

\section{Le call/cc et la logique classique}

\subsection{Continuations}

\thingy Dans un langage de programmation, la \defin{continuation} d'un
appel de fonction est l'état du calcul qui attend que cette fonction
renvoie une valeur. On peut la voir comme une copie de la pile d'appels
au moment de l'appel.

Certains langages permettent de capturer cette continuation, de la
stocker comme une valeur et de l'invoquer explicitement. Invoquer une
continuation avec une valeur $v$ restaure la pile capturée et fait
comme si l'appel concerné avait renvoyé $v$, même si ce calcul avait
déjà terminé. Une continuation est ainsi comparable à un
\texttt{goto} qui restaure en même temps la pile d'appels.

Les continuations peuvent remplacer certains mécanismes d'exception,
mais aussi reprendre un calcul interrompu. Elles permettent notamment
d'implémenter des générateurs et certaines formes de multitâche
coopératif. Leur coût d'implémentation porte essentiellement sur la
gestion de la pile.

\thingy La fonction \defin{call/cc}, abréviation de
\textit{call-with-current-continuation}, prend en argument une fonction
$g$, capture sa propre continuation et passe celle-ci à $g$.

Une continuation ne retourne jamais au point où elle est invoquée.
Elle peut donc être typée comme $\alpha\to\bot$, ou plus généralement
comme $\alpha\to\beta$ pour un type de retour arbitraire $\beta$.
La fonction $g$ attend une continuation de type $\alpha\to\beta$ et
renvoie une valeur de type $\alpha$. On obtient ainsi
\[
\texttt{call/cc}:
((\alpha\to\beta)\to\alpha)\to\alpha.
\]
Ce type correspond à la loi de Peirce
\[
((A\Rightarrow B)\Rightarrow A)\Rightarrow A,
\]
qui est une formule caractéristique de la logique classique.

La présence de \texttt{call/cc} transforme donc la logique du typage :
la correspondance de Curry--Howard ne donne plus seulement une logique
intuitionniste, mais une logique classique. Cette interprétation peut
être formalisée par le $\lambda\mu$-calcul.

\thingy Le tiers exclu
\[
A\lor\neg A
\]
peut alors être implémenté en capturant une continuation. L'idée est de
renvoyer provisoirement le cas $\neg A$. Si cette prétendue négation
est ensuite appliquée à une valeur de type $A$, la continuation
capturée est invoquée afin de revenir en arrière et de renvoyer cette
fois le cas $A$.

Cette possibilité montre pourquoi la logique classique est moins
constructive : une valeur somme annoncée comme appartenant à l'un des
deux cas n'est plus stable, puisqu'une continuation peut faire revenir
le calcul dans le passé et modifier le choix.

\subsection{Continuation Passing Style}

\thingy Même lorsqu'un langage ne fournit pas directement de
continuations de première classe, on peut réécrire les programmes en
\defin{Continuation Passing Style}, ou CPS. Au lieu de renvoyer une
valeur $v$, une fonction reçoit en argument une continuation $k$ et
appelle $k$ sur $v$. Tous les appels deviennent ainsi terminaux.

Fixons un type $Z$ de retour ultime et posons
\[
\mathop{\sim}P:=(P\Rightarrow Z),
\qquad
\mathop{\sim}\mathop{\sim}P
  :=((P\Rightarrow Z)\Rightarrow Z).
\]
Une valeur de type $P$ passée par continuation a donc pour type
$\mathop{\sim}\mathop{\sim}P$.

Pour les termes du $\lambda$-calcul non typé, la transformation
essentielle est donnée par
\[
v^{\mathrm{CPS}}:=\lambda k.kv,
\]
\[
(\lambda v.t)^{\mathrm{CPS}}
 :=\lambda k.\,k(\lambda v.t^{\mathrm{CPS}}),
\]
et
\[
(fx)^{\mathrm{CPS}}
 :=\lambda k.\,
 f^{\mathrm{CPS}}
 \bigl(\lambda f_0.\,
 x^{\mathrm{CPS}}
 \bigl(\lambda x_0.\,f_0x_0k\bigr)\bigr).
\]
Cette dernière formule évalue d'abord $f$, puis $x$, puis applique la
fonction obtenue à la valeur obtenue et à la continuation du calcul
complet.

\thingy Pour typer systématiquement cette transformation, on définit
d'abord $P^\diamond$ par
\[
A^\diamond=A
\]
pour une variable de type $A$, et
\[
(P\Rightarrow Q)^\diamond
 :=P^\diamond\Rightarrow
   \mathop{\sim}\mathop{\sim}Q^\diamond.
\]
Les produits, sommes, $\top$ et $\bot$ sont transformés composante par
composante. On pose ensuite
\[
P^{\mathrm{CPS}}
 :=\mathop{\sim}\mathop{\sim}P^\diamond.
\]
La transformation des termes est définie de manière à garantir
\[
\text{si }\vdash t:P,
\qquad
\text{alors }\vdash t^{\mathrm{CPS}}:P^{\mathrm{CPS}}.
\]

Lorsque $Z=\bot$, cette transformation correspond à une traduction de
la logique classique dans la logique intuitionniste par double
négation. Elle donne notamment
\[
\mathsf{CPC}\vdash P
\quad\Longleftrightarrow\quad
\mathsf{IPC}\vdash P^{\mathrm{CPS}}.
\]
\review{Les diapositives indiquent plusieurs variantes de la traduction
par double négation, notamment celles de Gödel--Gentzen et de Glivenko.
Il faudra décider dans une version ultérieure jusqu'où développer leur
comparaison.}

\section{Sémantique du calcul propositionnel intuitionniste}

\subsection{Sémantiques, correction et complétude}

\thingy Une logique $\mathsf{L}$ est définie syntaxiquement par des
formules, des axiomes et des règles de déduction. On note
$\mathsf{L}\vdash\varphi$ lorsque $\varphi$ est démontrable dans cette
logique.

Une \defin{sémantique} est une classe de modèles donnant un sens aux
symboles de la logique. On note
\[
\mathcal{M}\vDash\varphi
\]
lorsque la formule $\varphi$ est vraie dans le modèle $\mathcal{M}$.

Une sémantique est \defin{correcte} lorsque tout théorème est vrai dans
tout modèle :
\[
\mathsf{L}\vdash\varphi
\quad\Longrightarrow\quad
\mathsf{S}\vDash\varphi.
\]
Elle est \defin{complète} lorsque toute formule vraie dans tous les
modèles est démontrable :
\[
\mathsf{S}\vDash\varphi
\quad\Longrightarrow\quad
\mathsf{L}\vdash\varphi.
\]

La sémantique booléenne est correcte et complète pour la logique
classique, mais elle n'est que correcte pour la logique
intuitionniste. Il faut donc chercher des modèles plus riches pour
caractériser exactement cette dernière.

\subsection{Cadres de Kripke}

\begin{defn}
Un \defin{cadre de Kripke} est un ensemble partiellement ordonné
$(W,\leq)$. Ses éléments sont appelés des \defin{mondes}. Si
$w\leq w'$, on dit que $w'$ est accessible depuis $w$.

Une affectation de vérité sur ce cadre associe à chaque variable
propositionnelle une fonction
\[
p:W\to\{0,1\}
\]
croissante : si $p(w)=1$ et $w\leq w'$, alors $p(w')=1$. Cette
propriété est appelée permanence.

Un modèle de Kripke est un cadre muni d'une telle affectation pour
chaque variable propositionnelle.
\end{defn}

\thingy Les connecteurs sont interprétés dans un monde $w$ par
\begin{itemize}
\item $(q_1\mathbin{\dottedland}q_2)(w)=1$ si et seulement si
  $q_1(w)=q_2(w)=1$ ;
\item $(q_1\mathbin{\dottedlor}q_2)(w)=1$ si et seulement si
  $q_1(w)=1$ ou $q_2(w)=1$ ;
\item $(q_1\mathbin{\dottedlimp}q_2)(w)=1$ si et seulement si, pour
  tout $w'\geq w$, $q_1(w')=1$ implique $q_2(w')=1$ ;
\item $\dottedtop(w)=1$ et $\dottedbot(w)=0$.
\end{itemize}
En particulier,
\[
\dot{\neg}p
 :=p\mathbin{\dottedlimp}\dottedbot
\]
vaut dans $w$ lorsque $p$ est faux dans tout monde accessible depuis
$w$.

Ainsi, $A\Rightarrow B$ ne signifie pas simplement que $A$ est faux ou
que $B$ est vrai dans le monde actuel. Il signifie que, dans tout
monde futur où $A$ deviendra vrai, $B$ sera également vrai.

Une formule est valide dans un modèle lorsqu'elle vaut dans tous ses
mondes, et valide pour la sémantique de Kripke lorsqu'elle vaut dans
tout modèle de Kripke.

\begin{thm}[Kripke]
La sémantique de Kripke est correcte et complète pour le calcul
propositionnel intuitionniste.
\end{thm}

\review{La preuve du théorème de Kripke n'est pas donnée dans les
diapositives et n'est donc pas développée ici.}

\thingy La complétude exige de considérer tous les cadres. Un cadre
réduit à un singleton redonne la sémantique booléenne. Un cadre
totalement ordonné valide
\[
(A\Rightarrow B)\lor(B\Rightarrow A),
\]
qui n'est pourtant pas démontrable intuitionnistement en général.

Le cadre à deux mondes $u\leq v$ donne une logique à trois valeurs
$\{0,\mathsf{q},1\}$, où $\mathsf{q}$ signifie qu'une proposition est
fausse dans $u$ mais vraie dans $v$.

\subsection{Sémantique topologique}

\thingy Fixons un espace topologique $X$ et notons $\mathcal{O}(X)$ son
ensemble d'ouverts. On interprète les connecteurs sur les ouverts par
\[
U\mathbin{\dottedland}V:=U\cap V,
\qquad
U\mathbin{\dottedlor}V:=U\cup V,
\]
\[
\dottedtop:=X,
\qquad
\dottedbot:=\varnothing,
\]
et
\[
U\mathbin{\dottedlimp}V
 :=\operatorname{int}\bigl((X\setminus U)\cup V\bigr).
\]
Autrement dit, $U\mathbin{\dottedlimp}V$ est le plus grand ouvert $W$
tel que
\[
U\cap W\subseteq V.
\]
On obtient notamment
\[
\dot{\neg}U=\operatorname{int}(X\setminus U)
\]
et
\[
\dot{\neg}\dot{\neg}U
 =\operatorname{int}(\operatorname{adh}(U)).
\]

Une formule est interprétée en substituant des ouverts à ses variables.
Elle est valide si son interprétation est toujours l'ouvert $X$.

\begin{thm}[Tarski]
La sémantique des ouverts est correcte et complète pour le calcul
propositionnel intuitionniste. Pour la complétude, il suffit même de
considérer les espaces $\mathbb{R}^n$, pour $n\geq1$.
\end{thm}

Cette sémantique traduit l'intuition suivant laquelle la vérité est
locale : si une propriété est vraie en un point, elle doit rester vraie
dans un voisinage de ce point.

\thingy À titre d'exemple, dans $X=\mathbb{R}$, soit
$U={]0,1[}$. Alors
\[
\dot{\neg}U
 ={]-\infty,0[}\cup{]1,+\infty[},
\qquad
\dot{\neg}\dot{\neg}U=U.
\]
On vérifie également que
\[
(\dot{\neg}\dot{\neg}U\mathbin{\dottedlimp}U)=\mathbb{R}.
\]
En revanche,
\[
U\mathbin{\dottedlor}\dot{\neg}U
 =\mathbb{R}\setminus\{0,1\},
\]
ce qui fournit un contre-modèle topologique au tiers exclu.

\subsection{Réalisabilité propositionnelle}

\thingy Reprenons une numérotation
\[
\Phi_e:\mathbb{N}\dashrightarrow\mathbb{N}
\]
des fonctions générales récursives. On interprète maintenant une
formule par un ensemble d'entiers, ses réalisateurs.

Pour $P,Q\subseteq\mathbb{N}$, on pose
\[
P\mathbin{\dottedland}Q
 :=\{\langle m,n\rangle:m\in P,\ n\in Q\},
\]
\[
P\mathbin{\dottedlor}Q
 :=\{\langle0,m\rangle:m\in P\}
   \cup\{\langle1,n\rangle:n\in Q\},
\]
et
\[
P\mathbin{\dottedlimp}Q
 :=\{e\in\mathbb{N}:
     \forall m\in P,\;
     \Phi_e(m)\mathord{\downarrow}\in Q\}.
\]
Enfin,
\[
\dottedtop:=\mathbb{N},
\qquad
\dottedbot:=\varnothing.
\]
Une formule $\varphi(A_1,\ldots,A_r)$ est dite
\defin{réalisable} lorsqu'il existe un même entier $n$ appartenant à
\[
\bigcap_{P_1,\ldots,P_r\subseteq\mathbb{N}}
\dot{\varphi}(P_1,\ldots,P_r).
\]

\begin{thm}[Dragalin, Troelstra, Nelson]
La sémantique de la réalisabilité est correcte pour le calcul
propositionnel intuitionniste.
\end{thm}

Le réalisateur peut être extrait de la preuve : il s'agit d'une autre
forme de la correspondance de Curry--Howard.

Cette sémantique n'est toutefois pas complète. Il existe des formules
réalisables qui ne sont pas démontrables intuitionnistement.

\thingy La formule $A\land B\Rightarrow B\land A$ est réalisée par un
programme qui permute les composantes d'un couple. C'est exactement
l'algorithme correspondant au terme
\[
\lambda z.\langle\pi_2z,\pi_1z\rangle
\]
déjà obtenu par Curry--Howard.

En revanche, $A\lor\neg A$ ne peut pas être uniformément réalisé : il
faudrait produire le même entier qui, pour toute partie
$P\subseteq\mathbb{N}$, indique soit un élément de $P$, soit un
programme certifiant que $P$ est vide.

\subsection{Sémantique des problèmes finis}

\thingy Un \defin{problème fini} est un couple $(X,S)$ où $X$ est un
ensemble fini non vide de candidats et $S\subseteq X$ l'ensemble de ses
solutions.

On définit
\[
(X,S)\mathbin{\dottedland}(Y,T)
 :=(X\times Y,S\times T),
\]
\[
(X,S)\mathbin{\dottedlor}(Y,T)
 :=(X\uplus Y,S\uplus T),
\]
et
\[
(X,S)\mathbin{\dottedlimp}(Y,T)
 :=(Y^X,U),
\]
où
\[
U:=\{f:X\to Y:f(S)\subseteq T\}.
\]
Enfin,
\[
\dottedtop:=(\{\bullet\},\{\bullet\}),
\qquad
\dottedbot:=(\{\bullet\},\varnothing).
\]

Une formule est dite valide au sens de Medvedev lorsqu'il existe, pour
tous les ensembles de candidats choisis, un même candidat qui est une
solution du problème interprétant la formule, quels que soient les
sous-ensembles de solutions.

\begin{thm}[Medvedev]
La sémantique des problèmes finis est correcte pour le calcul
propositionnel intuitionniste.
\end{thm}

Elle n'est pas complète. Elle formalise néanmoins de manière précise
l'idée de Brouwer--Heyting--Kolmogorov selon laquelle une preuve d'une
implication est une fonction transformant les solutions du premier
problème en solutions du second.

\section{Inférence de type à la Hindley--Milner}

\subsection{Désannotation et problème de l'inférence}

\thingy Le $\lambda$CST présenté jusqu'ici porte toutes les annotations
de type. La \defin{désannotation} d'un terme consiste à les supprimer.
Par exemple,
\[
\lambda(x:\alpha).
\lambda(f:(\alpha\to\beta)\to\beta\to\gamma).
f(\lambda(h:\alpha\to\beta).hx)
\]
se désannote en
\[
\lambda xf.\,f(\lambda h.hx).
\]

Le problème de l'inférence de type est alors le suivant :
\begin{itemize}
\item décider si un terme non typé est la désannotation d'un terme du
  $\lambda$CST ;
\item si oui, retrouver ses annotations et son type ;
\item trouver même un type principal, c'est-à-dire le type le plus
  général dont tous les autres s'obtiennent par substitution.
\end{itemize}

L'algorithme de Hindley--Milner, également appelé algorithme de
Damas--Milner ou algorithme W, accomplit ces tâches. Il transforme les
variables de type opaques du $\lambda$CST en un véritable mécanisme de
polymorphisme.

\subsection{Collection et résolution des contraintes}

\thingy La première phase de l'algorithme consiste à collecter des
contraintes. On attribue une variable de type fraîche à chaque
variable et à chaque résultat intermédiaire.

Pour une variable $v$, on renvoie le type qui lui est associé dans le
contexte. Pour une abstraction $\lambda v.e$, on attribue à $v$ un
type frais $\eta$, on infère un type $\zeta$ pour $e$, puis on renvoie
$\eta\to\zeta$. Pour une application $fx$, si les types provisoires de
$f$ et $x$ sont $\zeta$ et $\xi$, on introduit un type frais $\eta$,
on ajoute la contrainte
\[
\zeta=(\xi\to\eta)
\]
et on attribue le type $\eta$ à l'application.

À titre d'exemple, pour
\[
\lambda xyz.xz(yz),
\]
on collecte
\[
x:\eta_1,\qquad y:\eta_2,\qquad z:\eta_3,
\]
\[
xz:\eta_4
\quad\text{avec}\quad
\eta_1=(\eta_3\to\eta_4),
\]
\[
yz:\eta_5
\quad\text{avec}\quad
\eta_2=(\eta_3\to\eta_5),
\]
et
\[
xz(yz):\eta_6
\quad\text{avec}\quad
\eta_4=(\eta_5\to\eta_6).
\]
Le type provisoire renvoyé est
\[
\eta_1\to\eta_2\to\eta_3\to\eta_6.
\]

\thingy La deuxième phase résout les contraintes par
\defin{unification}. Une substitution des variables de type est
cherchée de manière à rendre égaux les deux membres de chaque
contrainte.

L'algorithme procède ainsi :
\begin{itemize}
\item une contrainte déjà satisfaite est supprimée ;
\item si l'un des membres est une variable $\eta$ qui n'apparaît pas
  dans l'autre membre, on substitue cet autre membre à $\eta$ partout ;
\item si les deux membres sont des types fonctions, on les décompose
  composante par composante ;
\item dans tous les autres cas, l'unification échoue.
\end{itemize}
La condition interdisant de remplacer $\eta$ par un type contenant
$\eta$ est appelée \defin{test d'occurrence}. Elle interdit notamment
de résoudre
\[
\eta=(\eta\to\zeta).
\]

\thingy Reprenons
\[
t:=(\lambda xy.x)(\lambda x'y'.x').
\]
La première phase donne les variables de type
\[
x:\eta_1,\quad y:\eta_2,\quad
x':\eta_3,\quad y':\eta_4
\]
et la contrainte
\[
(\eta_1\to\eta_2\to\eta_1)
 =
((\eta_3\to\eta_4\to\eta_3)\to\eta_5).
\]
L'unification produit successivement
\[
\eta_1\mapsto(\eta_3\to\eta_4\to\eta_3)
\]
et
\[
\eta_5\mapsto(\eta_2\to\eta_3\to\eta_4\to\eta_3).
\]
On obtient finalement le typage
\[
\vdash
(\lambda(x:\eta_3\to\eta_4\to\eta_3).
 \lambda(y:\eta_2).x)
(\lambda(x':\eta_3).\lambda(y':\eta_4).x')
:
\eta_2\to\eta_3\to\eta_4\to\eta_3.
\]

\begin{thm}[Propriétés de l'algorithme de Hindley--Milner]
L'algorithme d'unification termine. Lorsqu'il renvoie une substitution,
celle-ci satisfait les contraintes. S'il existe une solution, il
renvoie une solution principale dont toute autre solution s'obtient
par substitution.
\end{thm}

\begin{proof}
La terminaison se montre par une double récurrence, d'abord sur le
nombre de variables de type restant à éliminer, puis sur la taille
totale des types présents dans les contraintes.

La correction se vérifie directement à chaque transformation du
système de contraintes. La propriété de principalité se montre par
induction sur les appels récursifs de l'algorithme : toute solution du
système initial doit effectuer les identifications imposées par
l'étape courante, puis se factorise par la solution renvoyée pour le
système simplifié.

\review{Les diapositives ne donnent que les idées de ces preuves. La
notion exacte de composition des substitutions et la mesure de
terminaison devront être détaillées dans une version ultérieure.}
\end{proof}

\thingy Le terme $\lambda x.xx$ n'est pas typable dans le
$\lambda$CST. La collecte des contraintes imposerait
\[
\eta_1=(\eta_1\to\eta_2),
\]
ce que le test d'occurrence interdit. Cette impossibilité est cohérente
avec la normalisation forte : si $\lambda x.xx$ était typable, le terme
\[
(\lambda x.xx)(\lambda x.xx)
\]
le serait aussi, alors qu'il se réduit indéfiniment à lui-même.

On peut étendre certains systèmes à des types récursifs sans
constructeur, en supprimant ce test. L'option
\texttt{ocaml -rectypes} permet par exemple d'accepter de tels types.
Cette extension doit cependant être utilisée avec précaution.

\subsection{Polymorphisme du \texttt{let} et restriction de valeur}

\thingy Dans un langage fonctionnel,
\[
\texttt{let }v=x\texttt{ in }t
\]
peut être vu comme un sucre syntaxique pour
\[
(\lambda v.t)x.
\]
Cette traduction pose toutefois un problème pour le polymorphisme. Si
une même valeur doit être utilisée à plusieurs types dans $t$, le
typage ordinaire d'une abstraction impose un type unique à $v$.

L'idée de Hindley--Milner est de typer d'abord $x$, puis de
\defin{généraliser} les variables de type qui n'apparaissent pas dans
le contexte extérieur. Chaque occurrence de $v$ dans $t$ reçoit
ensuite des variables de type fraîches obtenues par instanciation du
type généralisé.

Cette opération permet par exemple d'utiliser une même fonction
\texttt{twice} aussi bien sur des entiers que sur des booléens dans le
corps d'un même \texttt{let}.

\thingy Lorsque le langage possède des effets de bord et des références
mutables, une généralisation sans restriction serait dangereuse.
Considérons schématiquement
\[
\texttt{let r = ref (fun x -> x) in
  r := (fun x -> x+1); (!r) true}.
\]
Si le type de \texttt{r} était généralisé, la même cellule pourrait être
vue à la fois comme une référence vers une fonction
$\texttt{int}\to\texttt{int}$ et comme une référence vers une fonction
$\texttt{bool}\to\texttt{bool}$. Le programme appliquerait alors une
fonction entière à un booléen.

OCaml utilise donc la \defin{restriction de valeur} : le membre droit
d'un \texttt{let} n'est pleinement généralisé que lorsqu'il est une
valeur statique déterminable à la compilation. Dans un langage
purement fonctionnel comme Haskell, ce problème ne se pose pas de la
même manière, puisque les valeurs et les fonctions sont pures.

\chapter{Introduction aux quantificateurs}
\label{chapter-quantifiers}

{\footnotesize\noindent
\review{Premier brouillon produit avec l'aide d'un modèle d'IA : ce
chapitre doit être relu, complété et vérifié par l'enseignant avant toute
diffusion.}\par}

\section{Les quantificateurs et les systèmes logiques}

\subsection{Les quantificateurs comme conjonctions et disjonctions en
famille}

\thingy
Jusqu'à présent, nous avons surtout considéré le \defin{calcul
propositionnel}, c'est-à-dire une logique dans laquelle les formules
représentent des affirmations et dans laquelle on peut les combiner par
les connecteurs propositionnels
\[
\Rightarrow,\qquad \land,\qquad \lor,\qquad \top,\qquad \bot.
\]
Ces connecteurs permettent déjà d'exprimer beaucoup de raisonnements.
Ils ne permettent toutefois pas d'exprimer directement des affirmations
du genre « tout entier possède telle propriété » ou « il existe un objet
possédant telle propriété ».

\thingy
Pour cela, on introduit deux notations logiques fondamentales : les
\defin{quantificateurs} universel et existentiel, notés respectivement
$\forall$ et $\exists$.  Ils s'appliquent à une formule $P(v)$ qui
dépend d'une variable libre $v$, supposée ici de type $I$.  Le
quantificateur lie alors cette variable et forme une nouvelle formule
\[
\forall (v:I).\,P(v)
\qquad\hbox{ou}\qquad
\exists (v:I).\,P(v).
\]
Lorsque le type des individus est clair ou laissé implicite, on écrit
souvent plus brièvement
\[
\forall v.\,P(v)
\qquad\hbox{et}\qquad
\exists v.\,P(v).
\]

\thingy
Intuitivement, on peut voir les quantificateurs comme des « $\land$ et
$\lor$ en famille ».  Plus précisément, la formule $\forall v.\,P(v)$,
parfois notée $\bigwedge_v P(v)$, joue le rôle d'une conjonction portant
sur toute une famille de formules.  De même, $\exists v.\,P(v)$, parfois
notée $\bigvee_v P(v)$, joue le rôle d'une disjonction portant sur une
famille.

Cette analogie est la même que celle qui relie les produits et les
sommes indexés aux produits et sommes binaires :
\[
\forall v.\,P(v)
\quad\hbox{est à}\quad
P\land Q
\quad\hbox{ce que}\quad
\prod_i p_i
\quad\hbox{est à}\quad
p\times q,
\]
et
\[
\exists v.\,P(v)
\quad\hbox{est à}\quad
P\lor Q
\quad\hbox{ce que}\quad
\sum_i p_i
\quad\hbox{est à}\quad
p+q.
\]

\thingy
Il faut cependant faire attention au fait que cette présentation reste
pour l'instant informelle.  Il existe de nombreux systèmes logiques,
qui diffèrent notamment par les objets sur lesquels on a le droit de
quantifier.  Autrement dit, il faudra préciser ce que peut être la
variable $v$, quel est son domaine ou son type $I$, et comment sont
formés les termes qui représentent des individus de ce type.

\subsection{Interprétation de Brouwer--Heyting--Kolmogorov}

\thingy
L'interprétation de Brouwer--Heyting--Kolmogorov, ou interprétation
\defin{BHK}, donne une lecture constructive des formules logiques : une
formule est comprise à travers ce qui constituerait un témoignage, ou
une preuve, de cette formule.  Pour les connecteurs propositionnels,
nous avons déjà vu les idées suivantes :
\begin{itemize}
\item un témoignage de $P\land Q$ est la donnée d'un témoignage de $P$
  et d'un témoignage de $Q$ ;
\item un témoignage de $P\lor Q$ est la donnée d'un témoignage de $P$
  ou d'un témoignage de $Q$, ainsi que l'indication du cas choisi ;
\item un témoignage de $P\Rightarrow Q$ est un moyen de transformer
  tout témoignage de $P$ en un témoignage de $Q$ ;
\item un témoignage de $\top$ est trivial ;
\item il n'existe pas de témoignage de $\bot$.
\end{itemize}

\thingy
Les quantificateurs prolongent naturellement cette lecture.  Un
témoignage de
\[
\forall v.\,P(v)
\]
est un moyen qui, à partir d'un individu quelconque $x$, fournit un
témoignage de $P(x)$.  En revanche, un témoignage de
\[
\exists v.\,P(v)
\]
est la donnée d'un certain individu $t_0$, accompagnée d'un témoignage
de $P(t_0)$.

Autrement dit, pour démontrer universellement une propriété, il faut
savoir traiter un individu arbitraire ; pour démontrer existentiellement
une propriété, il faut exhiber un individu qui convient.

\subsection{Introduction et élimination de $\forall$ et $\exists$}

\thingy
Pour introduire une quantification universelle
$\forall (v:I).\,Q(v)$, on choisit une variable $v:I$ arbitraire.  Il
faut que cette variable soit \defin{fraîche}, c'est-à-dire qu'elle
n'apparaisse librement dans aucune des hypothèses actuellement en cours.
Si l'on démontre alors $Q(v)$ pour ce $v$ arbitraire, on peut conclure
\[
\forall (v:I).\,Q(v).
\]

Dans une rédaction mathématique, le schéma correspondant est le
suivant : « soit $v$ arbitraire ; montrons $Q(v)$ ; donc
$\forall (v:I).\,Q(v)$ ».  La condition de fraîcheur est essentielle :
si $v$ dépendait d'une hypothèse particulière, il ne serait plus
véritablement arbitraire.

\thingy
Pour éliminer une quantification universelle, on part d'une preuve de
$\forall (v:I).\,Q(v)$ et l'on choisit un terme $t$ de type $I$.  On
obtient alors la formule $Q[v\backslash t]$, c'est-à-dire la formule
$Q$ dans laquelle les occurrences libres de $v$ ont été remplacées par
$t$ :
\[
\inferrule
  {\Gamma\vdash \forall (v:I).\,Q
   \\
   \Gamma\vdash t:I}
  {\Gamma\vdash Q[v\backslash t]}.
\]
En termes plus ordinaires, si une propriété est vraie pour tout
individu, elle est vraie en particulier pour l'individu représenté par
$t$.

\thingy
Pour introduire une quantification existentielle
$\exists (v:I).\,Q(v)$, il suffit au contraire d'exhiber un terme $t$
de type $I$, puis de démontrer $Q[v\backslash t]$.  On peut alors
conclure
\[
\inferrule
  {\Gamma\vdash t:I
   \\
   \Gamma\vdash Q[v\backslash t]}
  {\Gamma\vdash \exists (v:I).\,Q}.
\]
Le terme $t$ est le témoin de l'existence affirmée.

\thingy
Enfin, pour éliminer une quantification existentielle, on suppose que
l'on dispose de $\exists (v:I).\,P(v)$ et que l'on cherche à démontrer
une conclusion $Q$.  On peut alors considérer un individu $v:I$
arbitraire satisfaisant $P(v)$, à condition que $v$ soit frais : il ne
doit apparaître librement ni dans les hypothèses déjà présentes, ni
dans la conclusion $Q$.  Si, sous les hypothèses supplémentaires
$v:I$ et $P(v)$, on arrive à démontrer $Q$, alors on peut conclure $Q$
à partir de l'existence initiale :
\[
\inferrule
  {\Gamma\vdash \exists (v:I).\,P
   \\
   \Gamma,\,v:I,\,P\vdash Q}
  {\Gamma\vdash Q}.
\]

La condition portant sur $Q$ exprime que la conclusion ne doit pas
dépendre du témoin particulier choisi.  On raisonne bien à partir de
l'existence d'un témoin quelconque, et non à partir d'un témoin que
l'on aurait le droit de retenir en dehors de la démonstration.

\thingy
On peut résumer les règles de déduction naturelle ainsi obtenues par le
tableau suivant.  Le contexte $\Gamma$ peut contenir à la fois des
formules, qui jouent le rôle d'hypothèses logiques, et des variables
d'individus avec leur type, comme $v:I$.
\[
\begin{array}{c|c|c}
 & \text{Introduction} & \text{Élimination}\\\hline
\forall &
\inferrule{\Gamma,\,v:I\vdash Q}
          {\Gamma\vdash\forall(v:I).\,Q}
\quad(v\ \text{frais})
&
\inferrule{\Gamma\vdash\forall(v:I).\,Q\\\Gamma\vdash t:I}
          {\Gamma\vdash Q[v\backslash t]}
\\[1em]\hline
\exists &
\inferrule{\Gamma\vdash t:I\\\Gamma\vdash Q[v\backslash t]}
          {\Gamma\vdash\exists(v:I).\,Q}
&
\inferrule{\Gamma\vdash\exists(v:I).\,P\\
           \Gamma,\,v:I,\,P\vdash Q}
          {\Gamma\vdash Q}
\quad(v\ \text{frais})
\end{array}
\]

\subsection{Preuves, termes d'individus et notations de Curry--Howard}

\thingy
Dans l'interprétation de Curry--Howard, la conjonction $P\land Q$
correspond à un type produit $\sigma\times\tau$.  Une preuve de la
conjonction contient donc une preuve de chacune des deux composantes,
de même qu'une valeur d'un produit contient une valeur de chacun des
deux types.

Par analogie, la quantification universelle $\forall v.\,P(v)$
correspond à un \defin{type produit en famille}
\[
\prod_v \sigma(v).
\]
Un élément d'un tel type est une fonction qui, pour chaque $v$, fournit
une valeur de type $\sigma(v)$.  Une preuve de
$\forall(v:I).\,P(v)$ se note donc naturellement sous la forme
\[
\lambda(v:I).\,(\cdots),
\]
où le corps de l'abstraction est une preuve de $P(v)$.

\thingy
Avec les notations des $\lambda$-termes, la règle d'introduction de
$\forall$ devient
\[
\inferrule
  {\Gamma,\,v:I\vdash s:Q}
  {\Gamma\vdash\lambda(v:I).\,s
   \;:\;\forall(v:I).\,Q}.
\]
La variable $v$ peut apparaître dans $Q$ et dans le terme $s$, mais elle
ne doit pas apparaître dans le contexte $\Gamma$.

L'élimination de $\forall$ est alors une application ordinaire :
\[
\inferrule
  {\Gamma\vdash f:\forall(v:I).\,Q
   \\
   \Gamma\vdash t:I}
  {\Gamma\vdash ft:Q[v\backslash t]}.
\]
Une preuve de $\forall(v:I).\,Q(v)$ se comporte donc comme une fonction
qui reçoit un individu $t$ de type $I$ et renvoie une preuve de $Q(t)$.

\thingy
De même, la disjonction $P\lor Q$ correspond à un type somme
$\sigma+\tau$.  La quantification existentielle $\exists v.\,P(v)$
correspond alors à un \defin{type somme en famille}
\[
\sum_v \sigma(v).
\]
Une valeur de ce type est la donnée d'un indice $t_0$, accompagnée
d'une valeur de type $\sigma(t_0)$.  Une preuve de
$\exists(v:I).\,P(v)$ prend ainsi la forme d'une paire
\[
\langle t_0,\cdots\rangle.
\]

\thingy
La règle d'introduction de $\exists$ peut se noter
\[
\inferrule
  {\Gamma\vdash t:I
   \\
   \Gamma\vdash z:Q[v\backslash t]}
  {\Gamma\vdash\langle t,z\rangle
   \;:\;\exists(v:I).\,Q}.
\]
Le premier composant est le témoin, et le second est la preuve que ce
témoin satisfait la propriété demandée.

Pour éliminer une preuve existentielle, on utilise un filtrage de motif
analogue à celui employé pour les types sommes :
\[
\inferrule
  {\Gamma\vdash z:\exists(v:I).\,P
   \\
   \Gamma,\,v:I,\,h:P\vdash s:Q}
  {\Gamma\vdash
   (\texttt{match~}z\texttt{~with~}\langle v,h\rangle\mapsto s)
   \;:\;Q}.
\]
Ici encore, $v$ peut apparaître dans $s$, mais ne doit apparaître ni
dans $\Gamma$ ni dans $Q$.

\thingy
Cette correspondance entre quantificateurs et produits ou sommes en
famille a toutefois une limite importante.  Il n'est pas toujours
possible, dans un système logique donné, d'extraire le témoin $t_0$
d'une preuve de $\exists v.\,P(v)$ afin de l'utiliser comme une donnée
ordinaire.  Dans certains systèmes, le témoin ne peut être exploité
qu'à l'intérieur d'une preuve logique.

\review{La portée exacte de cette restriction dépend du système de types
considéré.  Elle devra être vérifiée lors de l'harmonisation du présent
chapitre avec les chapitres consacrés à Curry--Howard et aux systèmes de
types dépendants.}

\subsection{Ce que l'on peut quantifier : systèmes de types et
$\lambda$-cube}

\thingy
Les règles précédentes ne suffisent pas encore à définir complètement
un système logique.  Il faut aussi préciser la formation des termes
d'individus, ainsi que la façon dont le monde des individus et le monde
des propositions interagissent.

On peut distinguer, selon les systèmes, deux mondes plus ou moins
séparés :
\begin{itemize}
\item le monde des termes, ou individus, et de leurs types ;
\item le monde logique, constitué des propositions et de leurs preuves.
\end{itemize}
Les règles de déduction naturelle présentées ci-dessus construisent des
démonstrations et appartiennent donc au monde logique.  Les prémisses de
la forme $\Gamma\vdash t:I$, qui assurent qu'un terme $t$ est bien un
individu de type $I$, relèvent du monde des termes.

\thingy
Dans Coq/Rocq, ces deux mondes sont séparés mais parallèles : les
propositions appartiennent à $\mathit{Prop}$, tandis que les types
d'individus appartiennent à $\mathit{Type}$.  D'autres systèmes adoptent
des séparations différentes, ou mélangent davantage ces deux niveaux.

Les questions suivantes permettent notamment de distinguer les
systèmes :
\begin{itemize}
\item peut-on quantifier sur les propositions, comme dans
  \[
  \forall(A:*).\,\bigl(A\Rightarrow A\bigr),
  \]
  avec la preuve
  $\lambda(A:*).\,\lambda(h:A).\,h$ ?
\item peut-on abstraire sur une proposition, et former par exemple
  \[
  \lambda(A:*).\,A
  \]
  de type $*\rightarrow *$ ?
\item peut-on quantifier sur des individus $x:I$ dans une famille de
  propositions $A:I\rightarrow *$, comme dans
  \[
  \forall(x:I).\,A(x)?
  \]
\end{itemize}

\thingy
Les réponses possibles à ces questions conduisent notamment aux huit
systèmes du \defin{$\lambda$-cube} de Barendregt.  Le but n'est pas ici
de les étudier en détail, mais de retenir que les règles de quantification
dépendent étroitement de ce que le système autorise comme objets, types,
propositions et abstractions.

\subsection{Quantification existentielle, types sommes et
imprédicativité}

\thingy
Considérons l'énoncé suivant, pour deux types $U$ et $V$ :
\[
\bigl(\forall(x:U).\,\exists(y:V).\,P(x,y)\bigr)
\Rightarrow
\bigl(\exists(f:U\Rightarrow V).\,\forall(x:U).\,P(x,f(x))\bigr).
\]
Cet énoncé est un analogue, dans le langage des types, de l'\defin{axiome
du choix}.  Il affirme que, si pour chaque $x:U$ on dispose d'un $y:V$
tel que $P(x,y)$, alors on peut réunir les témoins ainsi obtenus dans une
fonction $f:U\Rightarrow V$.

\thingy
Si $\forall$ et $\exists$ sont interprétés respectivement comme un
produit et une somme en famille, cet énoncé paraît naturel.  Une preuve
de
\[
\forall(x:U).\,\exists(y:V).\,P(x,y)
\]
est alors une fonction qui, à chaque $x$, associe une paire constituée
d'un $y$ et d'une preuve de $P(x,y)$.  On peut donc définir la fonction
qui extrait le premier composant de cette paire, et conserver dans une
seconde preuve le second composant.

En revanche, si $\exists$ est seulement un quantificateur logique, le
témoin fourni par une preuve existentielle ne peut pas nécessairement
être extrait comme une valeur utilisable dans le monde des programmes.
L'énoncé ci-dessus n'est alors pas nécessairement démontrable.

\thingy
Cette différence est particulièrement visible entre des systèmes tels
que Coq, où l'énoncé précédent n'est pas en général démontrable lorsque
$P:U\times V\to\mathit{Prop}$, et des systèmes à la Martin-Löf tels
qu'Agda, dans lesquels les preuves existentielles sont des paires
effectives et permettent cette construction.

\review{Cette comparaison entre Coq/Rocq et Agda reprend l'idée des
diapositives, mais devra être vérifiée avec soin au regard des univers,
des restrictions d'élimination de \texttt{Prop} et des bibliothèques
effectivement utilisées.}

\thingy
On appelle \defin{imprédicativité} la possibilité de définir une
proposition ou un type en quantifiant sur toutes les propositions ou
tous les types, y compris celui que l'on est en train de définir.  Il
s'agit donc d'une certaine forme de circularité.

Par exemple, en notant $*$ un « type des types » imprédicatif, on peut
considérer les expressions suivantes :
\[
A \;\cong\; \forall(Z:*).\,\bigl(Z\Rightarrow A\bigr),
\]
\[
A \;\cong\;
\forall(Z:*).\,\bigl((A\Rightarrow Z)\Rightarrow Z\bigr),
\]
\[
\bot\;\cong\;\forall(Z:*).\,Z,
\]
\[
A\land B\;\cong\;
\forall(Z:*).\,\bigl((A\Rightarrow B\Rightarrow Z)\Rightarrow Z\bigr),
\]
et
\[
A\lor B\;\cong\;
\forall(Z:*).\,\bigl((A\Rightarrow Z)\Rightarrow
(B\Rightarrow Z)\Rightarrow Z\bigr).
\]
Ces équivalences sont à comprendre ici comme des indications
informelles de codages possibles grâce à la quantification
imprédicative.

\thingy
L'imprédicativité est utile pour définir certaines constructions sur les
types.  Elle doit néanmoins être maniée avec prudence : sous des formes
trop générales, elle peut mener à des incohérences logiques, comme le
montre le paradoxe de Girard.

\review{Les codages imprédicatifs ci-dessus sont seulement esquissés dans
les diapositives.  Une passe ultérieure devra décider s'il convient de
donner les termes précis des deux sens de chaque isomorphisme, ou de
conserver cette présentation informelle.}

\section{Logique du premier ordre}

\subsection{Individus, prédicats et syntaxe des formules}

\thingy
La \defin{logique du premier ordre}, aussi appelée \defin{calcul des
prédicats} du premier ordre, est le système logique le plus simple qui
ajoute les quantificateurs au calcul propositionnel.  Les objets sur
lesquels il est permis de quantifier sont appelés des \defin{individus}.

Du point de vue des systèmes de types, cette logique est assez
rudimentaire : les individus forment un unique type implicite $I$,
spécialement réservé à ce rôle.  La logique pure du premier ordre ne
permet pas directement de construire des couples, des fonctions ou
d'autres objets à partir de ces individus.  Cette limitation est
compensée, dans les usages mathématiques ordinaires, par la théorie que
l'on exprime dans cette logique.

\thingy
La logique du premier ordre possède une importance mathématique
considérable.  La position classique est que les mathématiques peuvent
être formalisées dans la théorie des ensembles de Zermelo--Fraenkel avec
l'axiome du choix, notée $\mathsf{ZFC}$, elle-même formulée en logique
du premier ordre.  Les objets mathématiques élaborés, tels que couples,
fonctions ou structures, sont alors construits à l'intérieur de la
théorie des ensembles plutôt que fournis directement par la logique.

\thingy
Dans la logique du premier ordre pure, on dispose de deux sortes de
variables :
\begin{itemize}
\item des variables d'individus, notées par exemple $x$, $y$, $z$,
  en nombre illimité ;
\item pour chaque entier naturel $n$, des variables de prédicats
  $n$-aires, aussi appelées relations $n$-aires entre individus, notées
  par exemple $A^{(n)}$, $B^{(n)}$, $C^{(n)}$.
\end{itemize}
L'arité est généralement omise lorsqu'elle se lit immédiatement sur la
formule.

\thingy
Les formules sont définies inductivement.  Une formule est :
\begin{itemize}
\item une application $A^{(n)}(x_1,\ldots,x_n)$ d'une variable de
  prédicat $n$-aire à $n$ variables d'individus ;
\item une combinaison de deux formules $P$ et $Q$ par l'un des
  connecteurs $(P\Rightarrow Q)$, $(P\land Q)$ ou $(P\lor Q)$, ou bien
  l'une des constantes $\top$ et $\bot$ ;
\item une quantification $\forall x.\,P$ ou $\exists x.\,P$, qui lie la
  variable d'individu $x$ dans $P$.
\end{itemize}
Le type $I$ des individus reste implicite : les notations
$\forall x.\,P$ et $\exists x.\,P$ abrègent donc
$\forall(x:I).\,P$ et $\exists(x:I).\,P$.

\thingy
Il ne faut pas confondre les variables de prédicats avec les variables
d'individus.  Dans la logique du premier ordre, on peut quantifier sur
les individus, mais non sur les prédicats eux-mêmes : c'est précisément
le sens de l'expression « premier ordre ».

\subsection{Règles de déduction naturelle et conditions de fraîcheur}

\thingy
Les règles logiques de la logique du premier ordre reprennent les règles
du calcul propositionnel, auxquelles s'ajoutent les règles
d'introduction et d'élimination de $\forall$ et $\exists$.  Elles sont
les mêmes que celles données en
\ref{section-quantifiers-and-logical-systems}, mais la variable
quantifiée est désormais une variable d'individu $x$ du type implicite
$I$.

\review{Le label
\texttt{section-quantifiers-and-logical-systems} doit être ajouté à la
première section si l'on souhaite conserver ce renvoi.}

\thingy
Un point essentiel est que les prémisses de formation des termes,
comme
\[
\Gamma\vdash t:I,
\]
ne doivent pas être oubliées.  Elles assurent que le terme substitué
dans une formule universelle ou existentielle désigne effectivement un
individu.  Dans les arbres de preuve, ces prémisses sont parfois
omises pour alléger la présentation, mais elles restent nécessaires.

\thingy
Dans la variante la plus simple de la logique du premier ordre, les
seuls termes d'individus sont les variables d'individus.  Ainsi, la
seule façon d'obtenir une dérivation de $t:I$ est essentiellement
d'avoir déclaré $t:I$ dans le contexte.  La règle d'axiome doit donc
être comprise sous la forme
\[
\inferrule{}{\,\Gamma,\,x:I\vdash x:I}
\]
et, pour une formule atomique,
\[
\inferrule
  {\Gamma\vdash x_1:I\\\cdots\\\Gamma\vdash x_n:I}
  {\Gamma,\,A(x_1,\ldots,x_n)\vdash A(x_1,\ldots,x_n)}.
\]

\thingy
On ne suppose pas que le domaine des individus soit habité.  En
particulier, la formule
\[
\exists x.\,\top
\]
n'est pas démontrable sans hypothèse supplémentaire.  Pour la même
raison, on ne peut pas démontrer en général
\[
\bigl(\forall x.\,A(x)\bigr)\Rightarrow\bigl(\exists x.\,A(x)\bigr).
\]
Une preuve universelle ne fournit pas à elle seule un individu auquel
l'appliquer.

\subsection{Exemples de démonstrations et notation des preuves}

\thingy
Les formules propositionnelles sont des cas particuliers de formules du
premier ordre.  Il suffit, pour cela, de considérer les variables
propositionnelles comme des prédicats $0$-aires, aussi dits
\defin{nullaires}.  Ainsi,
\[
A\land B\Rightarrow B\land A
\]
est à la fois une formule propositionnelle et une formule du premier
ordre.

\thingy
Voici quelques exemples de formules du premier ordre :
\begin{itemize}
\item
  \[
  (\forall x.\,A(x))\land(\exists x.\,\top)
  \Rightarrow(\exists x.\,A(x));
  \]
\item
  \[
  (\forall x.\,\neg A(x))\Leftrightarrow(\neg\exists x.\,A(x));
  \]
\item
  \[
  (\exists x.\,\neg A(x))\Rightarrow(\neg\forall x.\,A(x));
  \]
\item
  \[
  (\exists x.\,A)\Leftrightarrow(\exists x.\,\top)\land A,
  \]
  où $A$ est un prédicat nullaire ;
\item
  \[
  (\forall x.\,A)\Leftrightarrow((\exists x.\,\top)\Rightarrow A),
  \]
  avec la même convention ;
\item
  \[
  (\exists x.\,\forall y.\,B(x,y))
  \Rightarrow
  (\forall y.\,\exists x.\,B(x,y)).
  \]
\end{itemize}

\thingy
Dans ces notes, nous suivons la convention selon laquelle les
quantificateurs $\forall$ et $\exists$ ont une priorité plus faible que
les connecteurs $\Rightarrow$, $\lor$, $\land$ et $\neg$.  Cette
convention n'est pas universelle : lorsque la moindre ambiguïté est
possible, il est préférable d'ajouter des parenthèses.

\thingy
Démontrons l'implication
\[
(\exists x.\,\forall y.\,B(x,y))
\Rightarrow
(\forall y.\,\exists x.\,B(x,y)).
\]
Supposons donc $\exists x.\,\forall y.\,B(x,y)$.  Soit $y'$ un individu
arbitraire.  Par élimination de l'existentielle, considérons un
individu $x$ tel que $\forall y.\,B(x,y)$.  Par élimination de
l'universelle, on obtient $B(x,y')$.  En introduisant l'existentielle
avec le témoin $x$, on en déduit $\exists x'.\,B(x',y')$.

Ce raisonnement ne dépend pas du choix de $y'$ ; on peut donc introduire
$\forall y'$.  De même, le témoin $x$ utilisé dans le raisonnement
intermédiaire était obtenu à partir de l'hypothèse existentielle ; on
peut donc conclure sans le conserver hors de cette élimination.  Enfin,
on décharge l'hypothèse initiale pour introduire l'implication.

\thingy
Avec les notations de Curry--Howard, la preuve précédente s'écrit
\[
\lambda(u:\exists x.\,\forall y.\,B(x,y)).\;
\lambda(y':I).\;
(\texttt{match~}u\texttt{~with~}\langle x,v\rangle\mapsto
\langle x,vy'\rangle).
\]
Le terme $u$ contient une paire existentielle : le filtrage de motif en
extrait un témoin $x$ et une preuve $v$ de $\forall y.\,B(x,y)$.  En
appliquant $v$ à $y'$, on obtient $vy':B(x,y')$, puis la paire
$\langle x,vy'\rangle$ constitue une preuve de
$\exists x'.\,B(x',y')$.

\subsection{Le monde des individus et le monde logique}

\thingy
La logique du premier ordre maintient une séparation particulièrement
nette entre le type implicite $I$ des individus et les relations
$A,B,C,\ldots$ qui servent à former des propositions.  Le type $I$ est
spécial : on ne l'écrit généralement pas, et la logique pure ne permet
pas de former directement des types tels que $I\times I$ ou $I\to I$.

\thingy
Les individus ne se mélangent donc pas directement aux preuves.  Les
variables $x$, $y$, $z$ appartiennent au monde des individus, tandis
que les formules et leurs démonstrations appartiennent au monde
logique.  Les prédicats permettent de parler logiquement des individus,
mais ils ne transforment pas pour autant les individus en propositions
ni les preuves en individus.

Cette organisation explique à la fois la simplicité de la logique du
premier ordre et ses limitations du point de vue du typ




%
%
%
\end{document}
